% Supplementary material for the ICDE 2027 EAB submission. Every number is a macro from ../macros_numbers.tex or a
% table generated by scripts/build_supplementary.py from reports/icde/tables; no number is typed by hand.
\documentclass[conference]{IEEEtran}
\IEEEoverridecommandlockouts
\usepackage[nocompress]{cite}
\usepackage{amsmath,amssymb,amsfonts}
\usepackage{graphicx}
\usepackage{booktabs}
\usepackage{multirow}
\usepackage{xcolor}
\usepackage[hidelinks]{hyperref}
\usepackage{url}
\graphicspath{{../figures/}}
\input{../macros_numbers.tex}

\begin{document}
\title{Supplementary Material for\\[2pt] [Experiment, Analysis, and Benchmark] CANTraj: Recovering Vehicle Trajectories from CAN Signals in Naturalistic Driving}
\author{\IEEEauthorblockN{Anonymous Author(s)}
\IEEEauthorblockA{Anonymous Institution(s)}}
\maketitle

\begin{abstract}
This document accompanies the paper with the material that supports its reproduction: the structure of the released
package and the commands that rebuild every table and figure (S1), the benchmark definitions in full (S2), the
reproduction details of the six published methods and of the benchmark baselines (S3), the complete result tables behind
the paper's findings, including the figures that the page limit left out (S4), the full protocol, prompts and
per-configuration results of the supplementary agent study (S5), and further recovered windows with their camera frames
(S6). Code: \url{https://github.com/WangYuHang-cmd/GPS--CAN}. Data, scores, configurations, logs and weights:
\url{https://huggingface.co/datasets/HenryYHW/CANTraj}.
\end{abstract}

\section{S1. Reproduction package}
The code repository holds the \texttt{cantraj} package (data layer, signal whitelist, causal replay, metrics and
scoring, baselines, map graphs, learned motion models, the agent study), the scripts that run every experiment and
build every table and figure, the frozen configuration (\texttt{configs/}: benchmark thresholds and manifest hashes,
signal registry with provenance classes, signal groups, protocols, windows, agent prompts and budgets), the
specification of benchmark v1 and the test suite (139 tests on synthetic data). The data package holds CANTraj
(\NReleaseDrives{} drives of \NReleaseDrivers{} drivers: re-decoded CAN event streams at logged times, the aligned GNSS
reference with the first and last two minutes of every drive cut, quality metadata and driver-disjoint splits), the
benchmark manifests (episodes, map-start rule, tile list with md5 and snapshot date, cohorts), every per-episode score
of every run, the run registry with code, scorer and tile hashes, the agent study's configuration and conversation
logs, the map tiles and the learned models' weights.

\paragraph{Rebuilding the tables} From the released per-episode scores alone: copy the package's \texttt{tables/} and
\texttt{reports/} into \texttt{reports/icde/}, then run \texttt{icde\_tables.py --extras a3 pairs topk signal cohort cost
robust robust\_map}, \texttt{icde\_area\_check.py}, \texttt{icde\_extra.py} (and \texttt{--what supp}),
\texttt{icde\_spy\_dewri\_diag.py}, \texttt{failure\_analysis.py}, \texttt{icde\_figures.py} and \texttt{icde\_macros.py}.
Every number of the paper is a macro in \texttt{macros\_numbers.tex} written by the last script.

\paragraph{Rerunning the methods} \texttt{run\_baseline.py} with the v1.1 flags (\texttt{--version v1.1 --manifest
episode\_manifest\_v1.parquet --eligible-only --tiles-dir maps/tiles\_v2 --map-episodes episode\_map\_v11.parquet}),
one chain script per study (\texttt{scripts/run\_icde\_*.sh}), then \texttt{score.py --run <id>}. Prediction files are
regenerated, not shipped. The agent study needs a served Qwen3-14B (revision pinned in the configuration) behind an
OpenAI-compatible endpoint; \texttt{scripts/h100\_setup.sh} installs and serves it, \texttt{run\_icde\_agent.sh} runs
the cache, the selector, the three seeds, the scoring and the tables.

\section{S2. Benchmark definitions}\label{sec:problem}
\paragraph{Episodes and eligibility} Windows of 60, 300 and 900\,s and whole drives tile each drive from its first
reference fix; a window starts at an anchor fix (speed $\ge 3$\,m/s, bearing accuracy within 5$^\circ$, horizontal
accuracy under 10\,m, CAN speed $\ge 2.5$\,m/s) and is eligible when it has at least 0.8 reference fixes per second, no
reference gap over 30\,s and a path of at least 100\,m. Drives with an unconfirmed vehicle platform and the controlled
supplement stay out of the main board. The map-start rule of v1.1 requires a drivable edge within 400\,m of the start
in the visible map; windows that fail it are reported separately and never counted against a method.
\paragraph{Splits} Driver-disjoint folds: train, validation, test and few-shot held-out (\NHeldoutDrives{} held-out drives
of \NHeldoutDrivers{} drivers). Parameters and models use train and validation only.
\paragraph{Signals and provenance} Signal sets S1 (speed), S2 (speed and pedals), S3 (speed and steering) and S4 (all);
groups G1 to G8 whitelist columns and static parameters; every signal carries a provenance class (bus-read, derived in
the vehicle software, logger sensor, GPS-informed), and class-F references are never in the CAN-only board.
\paragraph{Protocols} Online or offline access, map or no map, start heading given or withheld; the visible map is the
disc of radius $5\,\mathrm{km} + 40\,\mathrm{m/s} \cdot D$ (at most 60\,km) around the start, enforced per window.
\paragraph{Metrics} ADE and FDE at the reference seconds, coverage within 10 to 500\,m, time-free path recall and
precision at 25, 50 and 100\,m, route recovery (path recall at 50\,m $\ge 0.8$), terminal excursion and tail-lost
flags; an expected window without output is a failure (ADE $+\infty$, rates 0); medians carry driver-clustered
bootstrap intervals (1{,}000 resamples); paired differences are computed on common windows.

\section{S3. Methods}\label{sec:methods}
\input{../tables/methods.tex}
The six published methods were re-implemented from their papers and, where available, checked against the authors'
released code on the authors' own samples (OPD-II: \OpdSampleLogs{} of 59 logs within a median of at most
\OpdSampleMaxMedian{}\,m of the unmodified script; MacroTracking: the sample trace within a median of
\MacroSampleMedian{}\,m of the committed trajectory). Adaptations are named as such in Table~\ref{tab:methods}: Elastic
Pathing and Dewri et al.\ with the start heading (start edges within 45$^\circ$ of it), CAN-Trace and SPy with a 5\,km
square centred on the start, OPD-II with a strictly online read-out. Parameters left open by the papers were set on
training and validation drives only. The benchmark baselines are dead reckoning (nominal, calibrated, with an online
offset), a route prior and route retrieval, a map particle filter (2{,}000 particles, 10\,Hz ticks, heading and turn
noise of 15 and 20$^\circ$) and the offline map beam search (beam 200, ten candidates, heading and turn widths 15 and
20$^\circ$, U-turn weight 0.05, re-anchoring gain 0.3, length slack 2\,\% with a 10\,m floor, dead-end weight 0.1,
50\,m start radius, 10\,m merge bins, temperature 1), as given in Algorithm~1 of the paper. The learned motion models
(a GRU and a physics residual, three seeds) are trained on the training fold only with validation early stopping.

\section{S4. Complete results}
\begin{table*}[t]\centering\scriptsize\setlength{\tabcolsep}{3pt}
\caption{Complete board, 5 min windows, benchmark v1.1: expected windows $n$, failure share (no output), failure-aware median ADE in metres with the 95\,\% driver-clustered CI, and the route recovery rate (path recall at 50\,m $\ge 0.8$), test and few-shot drivers.}\label{tab:supp_main_w05}
\begin{tabular}{llrrrrrrrr}\toprule
method & protocol & $n$ & fail \% & median ADE m [CI] & route \% & $n$ & fail \% & median ADE m [CI] & route \% \\\midrule
hold start & O-N-Hp & 1{,}917 & 0.0 & 3{,}314 [2{,}113, 4{,}313] & 0 & 1{,}963 & 0.0 & 2{,}709 [2{,}232, 3{,}562] & 0 \\
straight line (speed) & O-N-Hp & 1{,}917 & 0.0 & 735 [644, 911] & 9 & 1{,}963 & 0.0 & 825 [781, 951] & 7 \\
start prior (speed, train stats) & O-N-Hp & 1{,}917 & 0.0 & 1{,}068 [920, 1{,}119] & 2 & 1{,}963 & 0.0 & 1{,}062 [946, 1{,}162] & 2 \\
route retrieval (speed, train routes) & O-N-Hp & 1{,}917 & 0.0 & 3{,}083 [2{,}002, 4{,}296] & 4 & 1{,}963 & 0.0 & 2{,}530 [2{,}016, 3{,}258] & 5 \\
steering DR, nominal & O-N-Hp & 1{,}917 & 0.0 & 921 [461, 1{,}742] & 4 & 1{,}963 & 0.0 & 1{,}669 [906, 2{,}472] & 2 \\
steering DR, online offset & O-N-Hp & 1{,}917 & 0.0 & 869 [611, 1{,}306] & 3 & 1{,}963 & 0.0 & 1{,}423 [964, 1{,}979] & 2 \\
steering DR, train calibrated & O-N-Hp & 1{,}917 & 0.0 & 1{,}312 [409, 4{,}189] & 4 & 1{,}963 & 0.0 & 1{,}651 [919, 2{,}430] & 2 \\
learned GRU (3 seeds) & O-N-Hp & 1{,}917 & 0.0 & 543 [226, 2{,}478] & 7 & 1{,}963 & 0.0 & 393 [305, 492] & 5 \\
learned physics residual (3 seeds) & O-N-Hp & 1{,}917 & 0.0 & 690 [310, 2{,}994] & 6 & 1{,}963 & 0.0 & 435 [366, 536] & 4 \\
map PF, speed only & O-M-Hp & 1{,}917 & 0.0 & 286 [50, 868] & 56 & 1{,}959 & 0.0 & 618 [186, 890] & 48 \\
map PF, steering & O-M-Hp & 1{,}917 & 0.0 & 43 [29, 59] & 87 & 1{,}959 & 0.0 & 27 [19, 47] & 90 \\
map PF, learned GRU motion (3 seeds) & O-M-Hp & 1{,}917 & 0.0 & 89 [42, 168] & 86 & 1{,}959 & 0.0 & 42 [29, 79] & 86 \\
map PF, learned residual motion (3 seeds) & O-M-Hp & 1{,}917 & 0.0 & 87 [45, 262] & 86 & 1{,}959 & 0.0 & 53 [33, 84] & 86 \\
MacroTracking & O-M-Hp & 1{,}917 & 0.0 & 84 [66, 210] & 62 & 1{,}959 & 0.0 & 52 [40, 143] & 65 \\
MacroTracking, train calibrated & O-M-Hp & 1{,}917 & 0.0 & 76 [61, 172] & 65 & 1{,}959 & 0.0 & 48 [38, 126] & 67 \\
OPD-II, online read-out & O-M-Hp & 1{,}917 & 0.0 & 42 [27, 69] & 79 & 1{,}959 & 0.0 & 24 [15, 54] & 78 \\
map beam, speed only & F-M-Hp & 1{,}917 & 0.0 & 165 [37, 633] & 54 & 1{,}959 & 0.0 & 343 [88, 637] & 47 \\
map beam, steering & F-M-Hp & 1{,}917 & 0.0 & 37 [24, 54] & 88 & 1{,}959 & 0.0 & 21 [17, 40] & 92 \\
map PF, steering, best-lineage read-out & F-M-Hp & 1{,}917 & 0.0 & 39 [28, 56] & 87 & 1{,}959 & 0.0 & 24 [17, 44] & 90 \\
Elastic Pathing (start heading used) & F-M-Hp & 1{,}917 & 4.6 & 87 [38, 324] & 58 & 1{,}959 & 4.0 & 80 [54, 165] & 59 \\
OPD-II (3 s delayed) & F-M-Hp & 1{,}917 & 0.0 & 42 [27, 68] & 79 & 1{,}959 & 0.0 & 24 [15, 54] & 78 \\
Dewri et al. (start heading used) & F-M-Hp & 1{,}917 & 0.9 & 81 [38, 215] & 62 & 1{,}959 & 2.8 & 89 [59, 157] & 59 \\
SPy, start known (start heading used) & F-M-Hp & 1{,}917 & 0.3 & 1{,}474 [1{,}283, 1{,}602] & 20 & 1{,}959 & 0.2 & 1{,}284 [1{,}090, 1{,}596] & 22 \\
map PF, steering & O-M-Hm & 1{,}917 & 0.0 & 138 [99, 175] & 79 & 1{,}959 & 0.0 & 88 [61, 124] & 82 \\
map PF, speed only & O-M-Hm & 1{,}917 & 0.0 & 1{,}773 [1{,}550, 1{,}987] & 34 & 1{,}959 & 0.0 & 1{,}868 [1{,}572, 2{,}097] & 29 \\
map beam, steering & F-M-Hm & 1{,}917 & 0.0 & 41 [26, 59] & 81 & 1{,}959 & 0.0 & 23 [17, 42] & 87 \\
map beam, speed only & F-M-Hm & 1{,}917 & 0.0 & 1{,}314 [866, 1{,}619] & 32 & 1{,}959 & 0.0 & 1{,}527 [1{,}294, 1{,}738] & 29 \\
Elastic Pathing & F-M-Hm & 1{,}917 & 3.2 & 129 [42, 691] & 53 & 1{,}959 & 2.1 & 177 [94, 344] & 50 \\
CAN-Trace (pedal subset), start known & F-M-Hm & 245 & 99.2 & fail [--, --] & 0 & 162 & 96.9 & fail [--, --] & 0 \\
CAN-Trace (pedal subset), 5 km area & F-M-Hm & 245 & 65.3 & fail [4{,}344, --] & 0 & 162 & 60.5 & fail [--, --] & 0 \\
Dewri et al. & F-M-Hm & 1{,}917 & 0.4 & 293 [81, 775] & 46 & 1{,}959 & 1.3 & 487 [272, 708] & 43 \\
SPy, start known & F-M-Hm & 1{,}917 & 0.1 & 2{,}198 [1{,}702, 2{,}636] & 11 & 1{,}959 & 0.1 & 1{,}810 [1{,}512, 2{,}211] & 13 \\
SPy, 5 km area & F-M-Hm & 1{,}917 & 11.8 & 3{,}929 [3{,}222, 4{,}406] & 0 & 1{,}959 & 5.5 & 3{,}400 [3{,}050, 3{,}953] & 0 \\
curvature DR (class F, not CAN-only) & O-N-Hp & 1{,}917 & 0.0 & 84 [74, 107] & 27 & 1{,}963 & 0.0 & 80 [70, 92] & 27 \\
map PF, curvature (class F) & O-M-Hp & 1{,}917 & 0.0 & 44 [29, 61] & 86 & 1{,}959 & 0.0 & 28 [19, 49] & 89 \\
\bottomrule\end{tabular}
\par\smallskip\raggedright Left block: test drivers; right block: few-shot drivers. Protocols O/F (online / offline), N/M (no map / map), Hp/Hm (start heading given / withheld).
\end{table*}

\begin{table*}[t]\centering\scriptsize\setlength{\tabcolsep}{3pt}
\caption{Complete board, 15 min windows, benchmark v1.1: expected windows $n$, failure share (no output), failure-aware median ADE in metres with the 95\,\% driver-clustered CI, and the route recovery rate (path recall at 50\,m $\ge 0.8$), test and few-shot drivers.}\label{tab:supp_main_w15}
\begin{tabular}{llrrrrrrrr}\toprule
method & protocol & $n$ & fail \% & median ADE m [CI] & route \% & $n$ & fail \% & median ADE m [CI] & route \% \\\midrule
hold start & O-N-Hp & 608 & 0.0 & 8{,}759 [5{,}761, 12{,}194] & 0 & 565 & 0.0 & 7{,}494 [5{,}973, 10{,}196] & 0 \\
straight line (speed) & O-N-Hp & 608 & 0.0 & 3{,}389 [3{,}124, 4{,}037] & 1 & 565 & 0.0 & 3{,}906 [3{,}630, 4{,}335] & 1 \\
start prior (speed, train stats) & O-N-Hp & 608 & 0.0 & 3{,}601 [3{,}128, 3{,}841] & 0 & 565 & 0.0 & 3{,}775 [3{,}545, 4{,}116] & 0 \\
route retrieval (speed, train routes) & O-N-Hp & 608 & 0.0 & 8{,}365 [5{,}617, 12{,}123] & 1 & 565 & 0.0 & 7{,}326 [5{,}648, 9{,}704] & 3 \\
steering DR, nominal & O-N-Hp & 608 & 0.0 & 5{,}135 [2{,}276, 7{,}673] & 1 & 565 & 0.0 & 7{,}174 [5{,}078, 9{,}162] & 0 \\
steering DR, online offset & O-N-Hp & 608 & 0.0 & 4{,}805 [3{,}893, 6{,}281] & 1 & 565 & 0.0 & 6{,}570 [5{,}331, 8{,}340] & 0 \\
steering DR, train calibrated & O-N-Hp & 608 & 0.0 & 6{,}375 [2{,}069, 12{,}288] & 1 & 565 & 0.0 & 7{,}138 [5{,}117, 9{,}181] & 0 \\
learned GRU (3 seeds) & O-N-Hp & 608 & 0.0 & 3{,}010 [1{,}193, 10{,}461] & 1 & 565 & 0.0 & 2{,}569 [1{,}923, 3{,}128] & 0 \\
learned physics residual (3 seeds) & O-N-Hp & 608 & 0.0 & 4{,}799 [1{,}785, 10{,}746] & 2 & 565 & 0.0 & 2{,}834 [2{,}292, 3{,}635] & 0 \\
map PF, speed only & O-M-Hp & 608 & 0.0 & 2{,}326 [364, 4{,}401] & 37 & 561 & 0.0 & 3{,}444 [2{,}231, 4{,}825] & 27 \\
map PF, steering & O-M-Hp & 608 & 0.0 & 178 [85, 373] & 69 & 561 & 0.0 & 127 [52, 224] & 73 \\
MacroTracking & O-M-Hp & 608 & 0.0 & 851 [293, 3{,}183] & 40 & 561 & 0.0 & 1{,}064 [245, 4{,}190] & 39 \\
MacroTracking, train calibrated & O-M-Hp & 608 & 0.0 & 779 [284, 2{,}875] & 39 & 561 & 0.0 & 1{,}110 [155, 3{,}912] & 40 \\
OPD-II, online read-out & O-M-Hp & 608 & 0.0 & 223 [100, 963] & 55 & 561 & 0.0 & 497 [94, 2{,}073] & 47 \\
map beam, speed only & F-M-Hp & 608 & 0.0 & 2{,}142 [136, 4{,}651] & 40 & 561 & 0.0 & 3{,}084 [1{,}498, 4{,}609] & 30 \\
map beam, steering & F-M-Hp & 608 & 0.0 & 149 [68, 331] & 73 & 561 & 0.0 & 106 [45, 174] & 76 \\
map PF, steering, best-lineage read-out & F-M-Hp & 608 & 0.0 & 165 [79, 394] & 68 & 561 & 0.0 & 125 [46, 215] & 72 \\
Elastic Pathing (start heading used) & F-M-Hp & 608 & 9.2 & 1{,}094 [124, 3{,}263] & 42 & 561 & 8.4 & 1{,}223 [928, 1{,}861] & 42 \\
OPD-II (3 s delayed) & F-M-Hp & 608 & 0.0 & 223 [100, 962] & 55 & 561 & 0.0 & 497 [94, 2{,}072] & 47 \\
Dewri et al. (start heading used) & F-M-Hp & 608 & 5.3 & 1{,}053 [142, 2{,}985] & 44 & 561 & 9.1 & 2{,}058 [1{,}159, 2{,}764] & 37 \\
SPy, start known (start heading used) & F-M-Hp & 608 & 0.5 & 5{,}396 [4{,}436, 6{,}077] & 6 & 561 & 0.5 & 4{,}819 [4{,}176, 6{,}440] & 6 \\
map PF, steering & O-M-Hm & 608 & 0.0 & 340 [201, 706] & 63 & 561 & 0.0 & 242 [118, 481] & 68 \\
map PF, speed only & O-M-Hm & 608 & 0.0 & 5{,}799 [4{,}754, 6{,}909] & 15 & 561 & 0.0 & 5{,}901 [4{,}652, 6{,}912] & 10 \\
map beam, steering & F-M-Hm & 608 & 0.0 & 155 [68, 376] & 70 & 561 & 0.0 & 110 [46, 184] & 74 \\
map beam, speed only & F-M-Hm & 608 & 0.0 & 5{,}099 [3{,}859, 6{,}250] & 25 & 561 & 0.0 & 5{,}419 [3{,}588, 6{,}052] & 21 \\
Elastic Pathing & F-M-Hm & 608 & 6.6 & 1{,}214 [135, 3{,}777] & 42 & 561 & 5.3 & 1{,}602 [1{,}001, 2{,}136] & 40 \\
CAN-Trace (pedal subset), 5 km area & F-M-Hm & 79 & 93.7 & fail [--, --] & 0 & 45 & 95.6 & fail [--, --] & 0 \\
Dewri et al. & F-M-Hm & 608 & 2.5 & 3{,}064 [1{,}642, 4{,}239] & 32 & 561 & 4.8 & 3{,}700 [2{,}838, 4{,}731] & 27 \\
SPy, start known & F-M-Hm & 608 & 0.2 & 6{,}199 [4{,}874, 7{,}741] & 4 & 561 & 0.4 & 5{,}928 [4{,}883, 7{,}689] & 4 \\
SPy, 5 km area & F-M-Hm & 608 & 6.6 & 9{,}871 [6{,}632, 12{,}216] & 0 & 561 & 4.8 & 8{,}349 [6{,}967, 10{,}942] & 0 \\
curvature DR (class F, not CAN-only) & O-N-Hp & 608 & 0.0 & 511 [445, 644] & 4 & 565 & 0.0 & 491 [424, 595] & 2 \\
map PF, curvature (class F) & O-M-Hp & 608 & 0.0 & 186 [86, 430] & 67 & 561 & 0.0 & 138 [61, 226] & 71 \\
\bottomrule\end{tabular}
\par\smallskip\raggedright Left block: test drivers; right block: few-shot drivers. Protocols O/F (online / offline), N/M (no map / map), Hp/Hm (start heading given / withheld).
\end{table*}

\begin{table*}[t]\centering\scriptsize\setlength{\tabcolsep}{3pt}
\caption{Complete board, whole drive windows, benchmark v1.1: expected windows $n$, failure share (no output), failure-aware median ADE in metres with the 95\,\% driver-clustered CI, and the route recovery rate (path recall at 50\,m $\ge 0.8$), test and few-shot drivers.}\label{tab:supp_main_full}
\begin{tabular}{llrrrrrrrr}\toprule
method & protocol & $n$ & fail \% & median ADE m [CI] & route \% & $n$ & fail \% & median ADE m [CI] & route \% \\\midrule
hold start & O-N-Hp & 282 & 0.0 & 5{,}235 [3{,}825, 7{,}854] & 0 & 460 & 0.0 & 3{,}322 [2{,}580, 4{,}870] & 0 \\
straight line (speed) & O-N-Hp & 282 & 0.0 & 6{,}551 [4{,}548, 9{,}073] & 2 & 460 & 0.0 & 3{,}327 [2{,}389, 4{,}744] & 8 \\
start prior (speed, train stats) & O-N-Hp & 282 & 0.0 & 5{,}771 [3{,}750, 7{,}481] & 2 & 460 & 0.0 & 2{,}912 [2{,}054, 3{,}825] & 6 \\
route retrieval (speed, train routes) & O-N-Hp & 282 & 0.0 & 5{,}586 [3{,}843, 7{,}854] & 1 & 460 & 0.0 & 3{,}236 [2{,}360, 4{,}800] & 3 \\
steering DR, nominal & O-N-Hp & 282 & 0.0 & 2{,}708 [1{,}468, 8{,}341] & 5 & 460 & 0.0 & 2{,}674 [1{,}251, 3{,}916] & 7 \\
steering DR, online offset & O-N-Hp & 282 & 0.0 & 5{,}004 [3{,}266, 6{,}809] & 4 & 460 & 0.0 & 3{,}118 [1{,}868, 4{,}449] & 8 \\
steering DR, train calibrated & O-N-Hp & 282 & 0.0 & 3{,}121 [1{,}917, 8{,}403] & 4 & 460 & 0.0 & 2{,}656 [1{,}498, 3{,}902] & 7 \\
learned GRU (3 seeds) & O-N-Hp & 282 & 0.0 & 1{,}704 [1{,}149, 4{,}563] & 6 & 460 & 0.0 & 887 [562, 1{,}547] & 11 \\
learned physics residual (3 seeds) & O-N-Hp & 282 & 0.0 & 1{,}944 [1{,}255, 5{,}372] & 6 & 460 & 0.0 & 1{,}100 [559, 1{,}838] & 10 \\
map PF, speed only & O-M-Hp & 282 & 0.0 & 5{,}827 [3{,}211, 7{,}147] & 8 & 454 & 0.0 & 2{,}908 [1{,}996, 3{,}946] & 13 \\
map PF, steering & O-M-Hp & 282 & 0.0 & 1{,}906 [799, 3{,}128] & 33 & 454 & 0.0 & 193 [91, 741] & 57 \\
MacroTracking & O-M-Hp & 282 & 0.0 & 1{,}649 [779, 4{,}479] & 33 & 454 & 0.0 & 227 [73, 900] & 46 \\
OPD-II, online read-out & O-M-Hp & 282 & 0.0 & 1{,}712 [913, 4{,}120] & 35 & 454 & 0.0 & 407 [52, 1{,}231] & 50 \\
map beam, speed only & F-M-Hp & 282 & 0.0 & 5{,}930 [4{,}134, 8{,}365] & 8 & 454 & 0.0 & 3{,}095 [2{,}044, 3{,}941] & 12 \\
map beam, steering & F-M-Hp & 282 & 0.0 & 1{,}179 [343, 1{,}729] & 41 & 454 & 0.0 & 94 [44, 239] & 61 \\
Elastic Pathing (start heading used) & F-M-Hp & 282 & 11.0 & 3{,}597 [1{,}952, 5{,}600] & 17 & 454 & 10.1 & 1{,}430 [957, 2{,}343] & 29 \\
OPD-II (3 s delayed) & F-M-Hp & 282 & 0.0 & 1{,}711 [912, 4{,}120] & 35 & 454 & 0.0 & 407 [53, 1{,}231] & 50 \\
Dewri et al. (start heading used) & F-M-Hp & 282 & 13.1 & 5{,}382 [2{,}969, 7{,}392] & 13 & 454 & 15.0 & 1{,}758 [1{,}119, 3{,}596] & 22 \\
SPy, start known (start heading used) & F-M-Hp & 282 & 1.4 & 4{,}752 [3{,}473, 7{,}599] & 4 & 454 & 0.9 & 2{,}651 [1{,}830, 4{,}240] & 6 \\
map PF, steering & O-M-Hm & 282 & 0.0 & 2{,}239 [1{,}076, 3{,}319] & 31 & 454 & 0.0 & 432 [230, 938] & 51 \\
map PF, speed only & O-M-Hm & 282 & 0.0 & 6{,}313 [3{,}834, 8{,}070] & 4 & 454 & 0.0 & 3{,}530 [2{,}312, 4{,}912] & 9 \\
map beam, steering & F-M-Hm & 282 & 0.0 & 1{,}244 [447, 2{,}085] & 39 & 454 & 0.0 & 124 [48, 510] & 58 \\
map beam, speed only & F-M-Hm & 282 & 0.0 & 6{,}958 [3{,}931, 9{,}102] & 6 & 454 & 0.0 & 3{,}605 [2{,}503, 4{,}779] & 9 \\
Elastic Pathing & F-M-Hm & 282 & 6.7 & 3{,}554 [2{,}359, 5{,}631] & 15 & 454 & 7.9 & 1{,}644 [1{,}056, 2{,}585] & 25 \\
Dewri et al. & F-M-Hm & 282 & 8.9 & 5{,}619 [3{,}273, 7{,}780] & 11 & 454 & 10.6 & 2{,}172 [1{,}429, 3{,}600] & 19 \\
SPy, start known & F-M-Hm & 282 & 1.8 & 4{,}932 [3{,}613, 7{,}779] & 2 & 454 & 1.1 & 3{,}254 [2{,}245, 4{,}387] & 3 \\
curvature DR (class F, not CAN-only) & O-N-Hp & 282 & 0.0 & 678 [498, 1{,}427] & 12 & 460 & 0.0 & 325 [204, 519] & 21 \\
\bottomrule\end{tabular}
\par\smallskip\raggedright Left block: test drivers; right block: few-shot drivers. Protocols O/F (online / offline), N/M (no map / map), Hp/Hm (start heading given / withheld).
\end{table*}

\begin{table*}[t]\centering\scriptsize\setlength{\tabcolsep}{3pt}
\caption{Same-cohort horizon study: one start per drive, scored 1, 5 and 15 minutes after the start (median ADE in metres / route recovery \%). Online methods on prefixes of their 15-minute runs; offline methods rerun on each prefix.}\label{tab:supp_cohort}
\begin{tabular}{llrrrrrr}\toprule
method & protocol & 1 min & 5 min & 15 min & 1 min & 5 min & 15 min \\\midrule
dr\_curvature & O-N-Hp & 13 / 81 & 101 / 18 & 569 / 4 & 13 / 91 & 105 / 19 & 577 / 3 \\
dr\_none & O-N-Hm & 382 / 4 & 2{,}150 / 0 & 7{,}693 / 0 & 439 / 7 & 2{,}749 / 0 & 8{,}806 / 0 \\
dr\_none & O-N-Hp & 103 / 28 & 1{,}364 / 1 & 5{,}317 / 0 & 82 / 31 & 1{,}141 / 5 & 5{,}131 / 1 \\
dr\_steer & O-N-Hm & 399 / 6 & 2{,}087 / 0 & 6{,}610 / 0 & 415 / 6 & 2{,}019 / 0 & 5{,}949 / 0 \\
dr\_steer & O-N-Hp & 24 / 60 & 500 / 4 & 3{,}308 / 0 & 40 / 48 & 941 / 3 & 5{,}066 / 0 \\
dr\_steer\_offonline & O-N-Hm & 371 / 5 & 1{,}904 / 0 & 5{,}618 / 0 & 425 / 5 & 2{,}108 / 0 & 6{,}367 / 0 \\
dr\_steer\_offonline & O-N-Hp & 36 / 46 & 1{,}080 / 2 & 5{,}066 / 0 & 37 / 45 & 1{,}174 / 1 & 5{,}548 / 0 \\
learned\_gru\_s0 & O-N-Hm & 385 / 7 & 2{,}189 / 0 & 7{,}032 / 0 & 426 / 6 & 2{,}331 / 0 & 6{,}733 / 0 \\
learned\_gru\_s0 & O-N-Hp & 20 / 66 & 282 / 10 & 1{,}483 / 2 & 25 / 61 & 301 / 4 & 1{,}989 / 0 \\
learned\_gru\_s1 & O-N-Hm & 388 / 7 & 2{,}255 / 0 & 6{,}927 / 0 & 424 / 7 & 2{,}361 / 1 & 6{,}740 / 0 \\
learned\_gru\_s1 & O-N-Hp & 18 / 67 & 234 / 7 & 1{,}239 / 1 & 19 / 66 & 289 / 6 & 1{,}896 / 0 \\
learned\_gru\_s2 & O-N-Hm & 384 / 6 & 2{,}207 / 0 & 7{,}107 / 0 & 425 / 7 & 2{,}260 / 0 & 6{,}884 / 0 \\
learned\_gru\_s2 & O-N-Hp & 19 / 68 & 300 / 7 & 1{,}438 / 2 & 21 / 67 & 308 / 5 & 2{,}008 / 0 \\
macrotracking & O-M-Hp & 20 / 83 & 62 / 59 & 1{,}351 / 35 & 20 / 93 & 46 / 62 & 953 / 37 \\
map\_pf\_curvature & O-M-Hm & 94 / 72 & 182 / 61 & 707 / 47 & 64 / 78 & 95 / 69 & 662 / 51 \\
map\_pf\_curvature & O-M-Hp & 19 / 81 & 65 / 63 & 708 / 47 & 16 / 93 & 35 / 75 & 232 / 56 \\
map\_pf\_none & O-M-Hm & 271 / 44 & 1{,}309 / 18 & 5{,}180 / 7 & 278 / 44 & 1{,}506 / 24 & 5{,}191 / 7 \\
map\_pf\_none & O-M-Hp & 113 / 56 & 1{,}145 / 21 & 4{,}720 / 12 & 92 / 60 & 1{,}037 / 31 & 4{,}383 / 12 \\
map\_pf\_steer & O-M-Hm & 90 / 74 & 205 / 60 & 815 / 47 & 51 / 82 & 75 / 72 & 483 / 54 \\
map\_pf\_steer & O-M-Hp & 19 / 82 & 56 / 64 & 672 / 49 & 15 / 93 & 31 / 77 & 249 / 56 \\
opd2\_online & O-M-Hp & 11 / 88 & 48 / 61 & 1{,}316 / 39 & 9 / 92 & 27 / 74 & 867 / 39 \\
elastic\_pathing & F-M-Hm & 70 / 49 & 703 / 34 & 3{,}335 / 20 & 29 / 55 & 316 / 41 & 2{,}379 / 28 \\
elastic\_pathing\_hp & F-M-Hp & 23 / 59 & 505 / 37 & 3{,}437 / 20 & 14 / 72 & 238 / 42 & 2{,}267 / 29 \\
map\_beam\_none & F-M-Hm & 167 / 31 & 1{,}274 / 14 & 5{,}352 / 11 & 212 / 34 & 1{,}358 / 27 & 4{,}763 / 11 \\
map\_beam\_none & F-M-Hp & 97 / 41 & 1{,}030 / 18 & 5{,}107 / 14 & 56 / 51 & 941 / 33 & 3{,}798 / 13 \\
map\_beam\_steer & F-M-Hm & 21 / 70 & 44 / 71 & 297 / 58 & 14 / 79 & 27 / 79 & 126 / 62 \\
map\_beam\_steer & F-M-Hp & 17 / 81 & 44 / 70 & 303 / 58 & 12 / 94 & 25 / 84 & 123 / 63 \\
map\_pf\_steer\_lineage & F-M-Hp & 17 / 81 & 53 / 66 & 675 / 47 & 13 / 93 & 27 / 77 & 241 / 55 \\
opd2 & F-M-Hp & 11 / 88 & 48 / 62 & 1{,}316 / 38 & 9 / 91 & 26 / 73 & 867 / 39 \\
can\_trace\_area & F-M-Hm & -- & -- & fail / 0 & -- & -- & fail / 0 \\
dewri & F-M-Hm & -- & -- & 3{,}312 / 20 & -- & -- & 3{,}700 / 18 \\
dewri\_hp & F-M-Hp & -- & -- & 3{,}069 / 25 & -- & -- & 3{,}562 / 22 \\
dr\_can\_yaw & O-N-Hm & -- & -- & fail / 0 & -- & -- & fail / 0 \\
dr\_can\_yaw & O-N-Hp & -- & -- & fail / 0 & -- & -- & fail / 0 \\
dr\_imu\_yaw\_rls & O-N-Hm & -- & -- & fail / 0 & -- & -- & fail / 0 \\
dr\_imu\_yaw\_rls & O-N-Hp & -- & -- & fail / 0 & -- & -- & fail / 0 \\
dr\_none\_pn05 & O-N-Hm & -- & -- & 7{,}695 / 0 & -- & -- & 8{,}806 / 0 \\
dr\_none\_pn05 & O-N-Hp & -- & -- & 5{,}318 / 0 & -- & -- & 5{,}145 / 1 \\
dr\_none\_pq1 & O-N-Hm & -- & -- & 7{,}694 / 0 & -- & -- & 8{,}816 / 0 \\
dr\_none\_pq1 & O-N-Hp & -- & -- & 5{,}306 / 0 & -- & -- & 5{,}116 / 1 \\
dr\_none\_s1hz & O-N-Hm & -- & -- & 7{,}693 / 0 & -- & -- & 8{,}806 / 0 \\
dr\_none\_s1hz & O-N-Hp & -- & -- & 5{,}317 / 0 & -- & -- & 5{,}131 / 1 \\
dr\_steer\_calib & O-N-Hm & -- & -- & 6{,}987 / 0 & -- & -- & 6{,}116 / 0 \\
dr\_steer\_calib & O-N-Hp & -- & -- & 3{,}138 / 1 & -- & -- & 5{,}052 / 0 \\
dr\_steer\_pn05 & O-N-Hm & -- & -- & 6{,}605 / 0 & -- & -- & 5{,}980 / 0 \\
dr\_steer\_pn05 & O-N-Hp & -- & -- & 3{,}286 / 1 & -- & -- & 5{,}047 / 0 \\
dr\_steer\_pq1 & O-N-Hm & -- & -- & 6{,}653 / 0 & -- & -- & 5{,}969 / 0 \\
dr\_steer\_pq1 & O-N-Hp & -- & -- & 3{,}473 / 0 & -- & -- & 5{,}069 / 0 \\
dr\_steer\_s1hz & O-N-Hm & -- & -- & 6{,}755 / 0 & -- & -- & 6{,}120 / 0 \\
dr\_steer\_s1hz & O-N-Hp & -- & -- & 3{,}278 / 2 & -- & -- & 4{,}939 / 0 \\
learned\_resid\_nogru\_s0 & O-N-Hm & -- & -- & 7{,}632 / 0 & -- & -- & 6{,}674 / 0 \\
learned\_resid\_nogru\_s0 & O-N-Hp & -- & -- & 3{,}898 / 0 & -- & -- & 4{,}718 / 0 \\
learned\_resid\_nogru\_s1 & O-N-Hm & -- & -- & 7{,}525 / 0 & -- & -- & 6{,}882 / 0 \\
learned\_resid\_nogru\_s1 & O-N-Hp & -- & -- & 3{,}792 / 0 & -- & -- & 5{,}024 / 0 \\
learned\_resid\_nogru\_s2 & O-N-Hm & -- & -- & 7{,}774 / 0 & -- & -- & 6{,}494 / 0 \\
learned\_resid\_nogru\_s2 & O-N-Hp & -- & -- & 3{,}719 / 0 & -- & -- & 4{,}664 / 0 \\
learned\_resid\_s0 & O-N-Hm & -- & -- & 6{,}780 / 0 & -- & -- & 6{,}760 / 0 \\
learned\_resid\_s0 & O-N-Hp & -- & -- & 2{,}253 / 2 & -- & -- & 2{,}285 / 0 \\
learned\_resid\_s1 & O-N-Hm & -- & -- & 7{,}090 / 0 & -- & -- & 6{,}557 / 0 \\
learned\_resid\_s1 & O-N-Hp & -- & -- & 1{,}910 / 2 & -- & -- & 2{,}050 / 0 \\
learned\_resid\_s2 & O-N-Hm & -- & -- & 6{,}743 / 0 & -- & -- & 6{,}903 / 0 \\
learned\_resid\_s2 & O-N-Hp & -- & -- & 1{,}803 / 2 & -- & -- & 2{,}582 / 0 \\
macrotracking\_calib & O-M-Hp & -- & -- & 1{,}076 / 35 & -- & -- & 985 / 39 \\
macrotracking\_wrapfix & O-M-Hp & -- & -- & 869 / 36 & -- & -- & 803 / 38 \\
opd2\_calib & F-M-Hp & -- & -- & 1{,}175 / 41 & -- & -- & 970 / 41 \\
opd2\_wrapfix & F-M-Hp & -- & -- & 1{,}316 / 38 & -- & -- & 867 / 39 \\
p0\_hold & O-N-Hm & -- & -- & 5{,}432 / 0 & -- & -- & 5{,}512 / 0 \\
p0\_hold & O-N-Hp & -- & -- & 5{,}432 / 0 & -- & -- & 5{,}512 / 0 \\
prior & O-N-Hm & -- & -- & 5{,}432 / 0 & -- & -- & 5{,}512 / 0 \\
prior & O-N-Hp & -- & -- & 4{,}391 / 0 & -- & -- & 4{,}328 / 0 \\
retrieval & O-N-Hm & -- & -- & 5{,}270 / 0 & -- & -- & 5{,}509 / 2 \\
retrieval & O-N-Hp & -- & -- & 5{,}270 / 0 & -- & -- & 5{,}509 / 2 \\
spy & F-M-Hm & -- & -- & 5{,}110 / 1 & -- & -- & 4{,}900 / 2 \\
spy\_area & F-M-Hm & -- & -- & 5{,}917 / 0 & -- & -- & 5{,}765 / 0 \\
spy\_hp & F-M-Hp & -- & -- & 4{,}761 / 2 & -- & -- & 4{,}592 / 2 \\
\bottomrule\end{tabular}
\par\smallskip\raggedright Left: test drives; right: few-shot drives.
\end{table*}

\begin{table*}[t]\centering\scriptsize\setlength{\tabcolsep}{3pt}
\caption{Signal study (pedal-valid five-minute windows of the evaluation drives): one learner, one decoder, per signal set and seed; failure-aware median ADE with 95\,\% driver-clustered CI, route recovery and 100\,m coverage.}\label{tab:supp_signal_sets}
\begin{tabular}{llrrrrrr}\toprule
setting & set & seed & $n$ & fail & median ADE m [CI] & route \% & c100 \% \\\midrule
no map & S1 & 0 & 407 & 0 & 1{,}079 [974, 1{,}209] & 1 & 19 \\
no map & S1 & 1 & 407 & 0 & 1{,}182 [1{,}005, 1{,}413] & 1 & 18 \\
no map & S1 & 2 & 407 & 0 & 1{,}128 [1{,}002, 1{,}324] & 1 & 19 \\
no map & S2c & 0 & 407 & 0 & 1{,}156 [1{,}006, 1{,}383] & 1 & 19 \\
no map & S2c & 1 & 407 & 0 & 1{,}156 [1{,}043, 1{,}323] & 0 & 18 \\
no map & S2c & 2 & 407 & 0 & 1{,}028 [854, 1{,}320] & 2 & 21 \\
no map & S3 & 0 & 407 & 0 & 960 [289, 1{,}437] & 7 & 29 \\
no map & S3 & 1 & 407 & 0 & 375 [226, 639] & 6 & 36 \\
no map & S3 & 2 & 407 & 0 & 376 [238, 728] & 6 & 35 \\
no map & S4c & 0 & 407 & 0 & 922 [312, 1{,}242] & 5 & 29 \\
no map & S4c & 1 & 407 & 0 & 401 [232, 759] & 8 & 38 \\
no map & S4c & 2 & 407 & 0 & 378 [274, 595] & 7 & 36 \\
no map & S2b & 0 & 407 & 0 & 1{,}119 [992, 1{,}369] & 1 & 19 \\
no map & S4b & 0 & 407 & 0 & 768 [264, 1{,}182] & 5 & 30 \\
map & S1 & 0 & 407 & 0 & 284 [169, 548] & 50 & 46 \\
map & S1 & 1 & 407 & 0 & 421 [332, 517] & 48 & 39 \\
map & S1 & 2 & 407 & 0 & 312 [201, 505] & 51 & 39 \\
map & S2c & 0 & 407 & 0 & 367 [290, 507] & 50 & 42 \\
map & S2c & 1 & 407 & 0 & 382 [294, 535] & 50 & 43 \\
map & S2c & 2 & 407 & 0 & 490 [394, 702] & 42 & 44 \\
map & S3 & 0 & 407 & 0 & 48 [33, 61] & 87 & 82 \\
map & S3 & 1 & 407 & 0 & 39 [29, 67] & 88 & 85 \\
map & S3 & 2 & 407 & 0 & 58 [25, 73] & 88 & 78 \\
map & S4c & 0 & 407 & 0 & 36 [28, 52] & 86 & 86 \\
map & S4c & 1 & 407 & 0 & 42 [36, 62] & 88 & 85 \\
map & S4c & 2 & 407 & 0 & 47 [34, 60] & 89 & 84 \\
map & S2b & 0 & 407 & 0 & 351 [246, 513] & 51 & 38 \\
map & S4b & 0 & 407 & 0 & 37 [28, 57] & 87 & 86 \\
\bottomrule\end{tabular}
\end{table*}

\begin{table*}[t]\centering\scriptsize\setlength{\tabcolsep}{3pt}
\caption{Signal study, paired gains of adding a signal (a $-$ b, metres; negative = better): median and mean paired differences with 95\,\% CIs and the share of windows that improve.}\label{tab:supp_signal_gains}
\begin{tabular}{lllllrrrr}\toprule
setting & added & a & b & seed & $n$ & median diff [CI] & mean diff [CI] & improved \% \\\midrule
no map & pedals | speed & S2c & S1 & 0 & 407 & 42.9 [14.7, 70.6] & 248 [107, 406] & 41 \\
no map & pedals | speed & S2c & S1 & 1 & 407 & -2.3 [-42.8, 54.0] & -159 [-257, 27] & 50 \\
no map & pedals | speed & S2c & S1 & 2 & 407 & -4.1 [-36.9, 20.1] & -81 [-233, 28] & 51 \\
no map & steering | speed & S3 & S1 & 0 & 407 & -134.5 [-515.2, 130.2] & 25 [-614, 358] & 56 \\
no map & steering | speed & S3 & S1 & 1 & 407 & -733.9 [-905.2, -352.2] & -853 [-1{,}110, -359] & 81 \\
no map & steering | speed & S3 & S1 & 2 & 407 & -530.6 [-747.3, -251.7] & -734 [-963, -288] & 75 \\
no map & pedals | speed+steering & S4c & S3 & 0 & 407 & -25.2 [-58.1, 14.7] & -85 [-135, 15] & 58 \\
no map & pedals | speed+steering & S4c & S3 & 1 & 407 & -3.7 [-22.3, 36.8] & -28 [-177, 110] & 52 \\
no map & pedals | speed+steering & S4c & S3 & 2 & 407 & -12.5 [-67.7, 45.8] & -58 [-219, 64] & 54 \\
no map & steering | speed+pedals & S4c & S2c & 0 & 407 & -258.1 [-658.4, -93.8] & -308 [-888, -21] & 63 \\
no map & steering | speed+pedals & S4c & S2c & 1 & 407 & -587.1 [-720.9, -419.4] & -722 [-898, -365] & 83 \\
no map & steering | speed+pedals & S4c & S2c & 2 & 407 & -405.2 [-652.0, -224.1] & -711 [-871, -413] & 75 \\
no map & pressed flags | speed & S2b & S1 & 0 & 407 & 76.2 [32.6, 102.2] & 162 [83, 279] & 40 \\
no map & pressed flags | speed+steering & S4b & S3 & 0 & 407 & -93.1 [-271.2, 5.0] & -218 [-325, -35] & 70 \\
map & pedals | speed & S2c & S1 & 0 & 407 & 0.9 [-1.6, 7.9] & 13 [-22, 40] & 48 \\
map & pedals | speed & S2c & S1 & 1 & 407 & -4.8 [-9.4, 0.6] & -2 [-48, 104] & 56 \\
map & pedals | speed & S2c & S1 & 2 & 407 & 6.1 [-6.6, 19.7] & 150 [-2, 235] & 44 \\
map & steering | speed & S3 & S1 & 0 & 407 & -119.5 [-305.3, -80.9] & -698 [-893, -581] & 83 \\
map & steering | speed & S3 & S1 & 1 & 407 & -233.7 [-363.5, -149.9] & -727 [-856, -582] & 87 \\
map & steering | speed & S3 & S1 & 2 & 407 & -147.0 [-343.4, -103.8] & -690 [-895, -582] & 88 \\
map & pedals | speed+steering & S4c & S3 & 0 & 407 & -6.4 [-10.9, -0.3] & -13 [-29, 20] & 72 \\
map & pedals | speed+steering & S4c & S3 & 1 & 407 & 0.9 [-1.2, 3.7] & 16 [-14, 55] & 44 \\
map & pedals | speed+steering & S4c & S3 & 2 & 407 & -3.7 [-9.0, 3.8] & -14 [-40, 8] & 63 \\
map & steering | speed+pedals & S4c & S2c & 0 & 407 & -191.3 [-320.7, -144.2] & -724 [-890, -583] & 90 \\
map & steering | speed+pedals & S4c & S2c & 1 & 407 & -231.4 [-348.8, -127.1] & -709 [-894, -570] & 84 \\
map & steering | speed+pedals & S4c & S2c & 2 & 407 & -359.2 [-461.9, -207.2] & -855 [-1{,}030, -621] & 85 \\
map & pressed flags | speed & S2b & S1 & 0 & 407 & 5.2 [0.8, 10.7] & 14 [-20, 77] & 41 \\
map & pressed flags | speed+steering & S4b & S3 & 0 & 407 & -3.6 [-12.2, 0.4] & -30 [-59, 2] & 67 \\
\bottomrule\end{tabular}
\end{table*}

\begin{table*}[t]\centering\scriptsize\setlength{\tabcolsep}{3pt}
\caption{Candidate sets on five-minute windows: whole-path oracle median ADE among the first $K$ candidates, share of windows whose best candidate is ranked first, and the median number of candidates.}\label{tab:supp_topk}
\begin{tabular}{lllrrrrrrrrr}\toprule
method & protocol & fold & $n$ & out & single & $K$=1 & 3 & 5 & 10 & best 1st \% & cands \\\midrule
map\_beam\_steer & F-M-Hp & test & 1{,}917 & 1{,}917 & 31 & 37 & 30 & 28 & 27 & 58 & 10 \\
map\_beam\_steer & F-M-Hp & few-shot & 1{,}959 & 1{,}959 & 15 & 21 & 18 & 18 & 17 & 61 & 10 \\
map\_beam\_steer & F-M-Hm & test & 1{,}917 & 1{,}917 & 2 & 41 & 32 & 29 & 27 & 54 & 10 \\
map\_beam\_steer & F-M-Hm & few-shot & 1{,}959 & 1{,}959 & 0 & 23 & 19 & 18 & 17 & 58 & 10 \\
map\_beam\_none & F-M-Hp & test & 1{,}917 & 1{,}917 & 31 & 165 & 85 & 73 & 60 & 35 & 10 \\
map\_beam\_none & F-M-Hp & few-shot & 1{,}959 & 1{,}959 & 18 & 343 & 89 & 73 & 56 & 28 & 10 \\
map\_beam\_none & F-M-Hm & test & 1{,}917 & 1{,}917 & 2 & 1{,}314 & 210 & 94 & 72 & 21 & 10 \\
map\_beam\_none & F-M-Hm & few-shot & 1{,}959 & 1{,}959 & 0 & 1{,}527 & 461 & 139 & 71 & 17 & 10 \\
elastic\_pathing\_hp & F-M-Hp & test & 1{,}917 & 1{,}828 & 1{,}402 & 87 & 77 & 76 & 75 & 84 & 1 \\
elastic\_pathing\_hp & F-M-Hp & few-shot & 1{,}959 & 1{,}881 & 1{,}488 & 80 & 66 & 65 & 64 & 85 & 1 \\
elastic\_pathing & F-M-Hm & test & 1{,}917 & 1{,}855 & 1{,}484 & 129 & 103 & 98 & 96 & 87 & 1 \\
elastic\_pathing & F-M-Hm & few-shot & 1{,}959 & 1{,}918 & 1{,}577 & 177 & 114 & 111 & 111 & 87 & 1 \\
can\_trace & F-M-Hm & test & 245 & 2 & 0 & fail & fail & fail & fail & 50 & 2 \\
can\_trace & F-M-Hm & few-shot & 162 & 5 & 2 & fail & fail & fail & fail & 80 & 2 \\
can\_trace\_area & F-M-Hm & test & 245 & 85 & 14 & fail & fail & fail & fail & 36 & 6 \\
can\_trace\_area & F-M-Hm & few-shot & 162 & 64 & 5 & fail & fail & fail & fail & 27 & 6 \\
dewri\_hp & F-M-Hp & test & 1{,}917 & 1{,}899 & 347 & 81 & 58 & 51 & 45 & 50 & 9 \\
dewri\_hp & F-M-Hp & few-shot & 1{,}959 & 1{,}905 & 311 & 89 & 49 & 41 & 35 & 47 & 6 \\
dewri & F-M-Hm & test & 1{,}917 & 1{,}909 & 40 & 293 & 70 & 60 & 49 & 39 & 10 \\
dewri & F-M-Hm & few-shot & 1{,}959 & 1{,}933 & 40 & 487 & 65 & 48 & 38 & 34 & 10 \\
\bottomrule\end{tabular}
\end{table*}

\begin{figure*}[t]\centering
\includegraphics[width=0.8\textwidth]{fig_topk_w05.pdf}
\caption{Best of the first $K$ candidates (whole-path oracle median ADE and per-second union coverage within 100\,m) for
the offline beam with and without steering and for Elastic Pathing and Dewri et al., five-minute windows.}
\end{figure*}
\begin{table*}[t]\centering\scriptsize\setlength{\tabcolsep}{3pt}
\caption{Fixed map attacks on degraded inputs (five-minute windows, grid fixed before the first run): median ADE and route recovery on the clean and the degraded signal, and the paired change of the route recovery rate in percentage points with 95\,\% driver-clustered CI.}\label{tab:supp_robust_map}
\begin{tabular}{llllrrrrrr}\toprule
method & protocol & degradation & fold & $n$ & med. clean & med. degraded & route clean \% & route degraded \% & $\Delta$ route pp [CI] \\\midrule
map PF, steering & O-M-Hp & steering quantised to 2 deg & test & 1{,}917 & 43 & 43 & 87 & 87 & +0.0 [-0.2, +0.3] \\
map PF, steering & O-M-Hp & steering quantised to 2 deg & few-shot & 1{,}959 & 27 & 27 & 90 & 90 & -0.2 [-0.6, +0.0] \\
map PF, steering & O-M-Hp & steering quantised to 5 deg & test & 1{,}917 & 43 & 42 & 87 & 87 & +0.2 [-0.2, +0.7] \\
map PF, steering & O-M-Hp & steering quantised to 5 deg & few-shot & 1{,}959 & 27 & 26 & 90 & 90 & -0.3 [-0.8, +0.1] \\
map PF, steering & O-M-Hp & steering quantised to 10 deg & test & 1{,}917 & 43 & 42 & 87 & 87 & +0.4 [-0.1, +0.9] \\
map PF, steering & O-M-Hp & steering quantised to 10 deg & few-shot & 1{,}959 & 27 & 27 & 90 & 90 & -0.6 [-1.3, -0.0] \\
map PF, steering & O-M-Hp & steering quantised to 20 deg & test & 1{,}917 & 43 & 44 & 87 & 86 & -1.4 [-2.5, -0.4] \\
map PF, steering & O-M-Hp & steering quantised to 20 deg & few-shot & 1{,}959 & 27 & 28 & 90 & 89 & -1.1 [-2.3, -0.4] \\
map PF, steering & O-M-Hp & steering at 1 Hz & test & 1{,}917 & 43 & 42 & 87 & 88 & +0.9 [+0.4, +1.8] \\
map PF, steering & O-M-Hp & steering at 1 Hz & few-shot & 1{,}959 & 27 & 27 & 90 & 91 & +0.8 [+0.3, +1.3] \\
map PF, steering & O-M-Hp & steering at 0.2 Hz & test & 1{,}917 & 43 & 49 & 87 & 84 & -2.7 [-5.8, -1.4] \\
map PF, steering & O-M-Hp & steering at 0.2 Hz & few-shot & 1{,}959 & 27 & 38 & 90 & 85 & -5.6 [-7.8, -3.8] \\
map PF, steering & O-M-Hp & steering offset noise 2 deg & test & 1{,}917 & 43 & 42 & 87 & 87 & -0.1 [-0.8, +0.3] \\
map PF, steering & O-M-Hp & steering offset noise 2 deg & few-shot & 1{,}959 & 27 & 27 & 90 & 90 & -0.3 [-0.8, +0.2] \\
map PF, steering & O-M-Hp & steering offset noise 5 deg & test & 1{,}917 & 43 & 43 & 87 & 87 & -0.5 [-1.2, +0.0] \\
map PF, steering & O-M-Hp & steering offset noise 5 deg & few-shot & 1{,}959 & 27 & 27 & 90 & 89 & -0.8 [-1.7, -0.1] \\
map PF, steering & O-M-Hp & steering delayed 0.5 s & test & 1{,}917 & 43 & 42 & 87 & 88 & +1.1 [+0.4, +2.1] \\
map PF, steering & O-M-Hp & steering delayed 0.5 s & few-shot & 1{,}959 & 27 & 27 & 90 & 91 & +0.6 [-0.3, +1.2] \\
map PF, steering & O-M-Hp & speed quantised to 1 m/s & test & 1{,}917 & 43 & 42 & 87 & 87 & +0.1 [-0.5, +0.5] \\
map PF, steering & O-M-Hp & speed quantised to 1 m/s & few-shot & 1{,}959 & 27 & 31 & 90 & 89 & -0.7 [-1.4, +0.5] \\
map PF, steering & O-M-Hp & speed quantised to 2 m/s & test & 1{,}917 & 43 & 41 & 87 & 87 & -0.2 [-0.7, +0.6] \\
map PF, steering & O-M-Hp & speed quantised to 2 m/s & few-shot & 1{,}959 & 27 & 42 & 90 & 88 & -2.6 [-4.1, -0.7] \\
map PF, steering & O-M-Hp & speed quantised to 5 m/s & test & 1{,}917 & 43 & 122 & 87 & 83 & -4.4 [-7.3, -2.4] \\
map PF, steering & O-M-Hp & speed quantised to 5 m/s & few-shot & 1{,}959 & 27 & 68 & 90 & 83 & -7.6 [-9.7, -5.0] \\
map PF, steering & O-M-Hp & speed at 1 Hz & test & 1{,}917 & 43 & 44 & 87 & 86 & -0.9 [-1.9, -0.1] \\
map PF, steering & O-M-Hp & speed at 1 Hz & few-shot & 1{,}959 & 27 & 29 & 90 & 89 & -1.3 [-1.8, -0.5] \\
map beam, steering & F-M-Hp & steering quantised to 2 deg & test & 1{,}917 & 37 & 37 & 88 & 88 & +0.1 [+0.0, +0.2] \\
map beam, steering & F-M-Hp & steering quantised to 2 deg & few-shot & 1{,}959 & 21 & 21 & 92 & 92 & -0.1 [-0.2, +0.0] \\
map beam, steering & F-M-Hp & steering quantised to 5 deg & test & 1{,}917 & 37 & 36 & 88 & 88 & +0.2 [-0.1, +0.8] \\
map beam, steering & F-M-Hp & steering quantised to 5 deg & few-shot & 1{,}959 & 21 & 22 & 92 & 92 & -0.1 [-0.4, +0.2] \\
map beam, steering & F-M-Hp & steering quantised to 10 deg & test & 1{,}917 & 37 & 37 & 88 & 88 & +0.0 [-0.3, +0.4] \\
map beam, steering & F-M-Hp & steering quantised to 10 deg & few-shot & 1{,}959 & 21 & 22 & 92 & 92 & -0.3 [-0.7, +0.1] \\
map beam, steering & F-M-Hp & steering quantised to 20 deg & test & 1{,}917 & 37 & 39 & 88 & 86 & -1.4 [-2.4, -0.3] \\
map beam, steering & F-M-Hp & steering quantised to 20 deg & few-shot & 1{,}959 & 21 & 25 & 92 & 91 & -1.1 [-2.2, -0.4] \\
map beam, steering & F-M-Hp & steering at 1 Hz & test & 1{,}917 & 37 & 37 & 88 & 89 & +1.0 [+0.2, +2.0] \\
map beam, steering & F-M-Hp & steering at 1 Hz & few-shot & 1{,}959 & 21 & 22 & 92 & 93 & +1.3 [+0.8, +2.0] \\
map beam, steering & F-M-Hp & steering at 0.2 Hz & test & 1{,}917 & 37 & 46 & 88 & 89 & +1.0 [-0.8, +3.0] \\
map beam, steering & F-M-Hp & steering at 0.2 Hz & few-shot & 1{,}959 & 21 & 36 & 92 & 90 & -1.8 [-3.4, -0.1] \\
map beam, steering & F-M-Hp & steering offset noise 2 deg & test & 1{,}917 & 37 & 37 & 88 & 87 & -0.2 [-0.4, +0.0] \\
map beam, steering & F-M-Hp & steering offset noise 2 deg & few-shot & 1{,}959 & 21 & 22 & 92 & 92 & -0.2 [-0.4, +0.2] \\
map beam, steering & F-M-Hp & steering offset noise 5 deg & test & 1{,}917 & 37 & 37 & 88 & 87 & -0.6 [-1.4, -0.1] \\
map beam, steering & F-M-Hp & steering offset noise 5 deg & few-shot & 1{,}959 & 21 & 23 & 92 & 92 & -0.1 [-0.5, +0.5] \\
map beam, steering & F-M-Hp & steering delayed 0.5 s & test & 1{,}917 & 37 & 37 & 88 & 89 & +1.3 [+0.5, +2.5] \\
map beam, steering & F-M-Hp & steering delayed 0.5 s & few-shot & 1{,}959 & 21 & 22 & 92 & 94 & +1.6 [+0.8, +2.4] \\
map beam, steering & F-M-Hp & speed quantised to 1 m/s & test & 1{,}917 & 37 & 36 & 88 & 88 & +0.0 [-0.7, +0.6] \\
map beam, steering & F-M-Hp & speed quantised to 1 m/s & few-shot & 1{,}959 & 21 & 26 & 92 & 91 & -0.7 [-1.2, +0.3] \\
map beam, steering & F-M-Hp & speed quantised to 2 m/s & test & 1{,}917 & 37 & 37 & 88 & 87 & -0.1 [-0.7, +0.8] \\
map beam, steering & F-M-Hp & speed quantised to 2 m/s & few-shot & 1{,}959 & 21 & 34 & 92 & 90 & -2.0 [-3.6, +0.0] \\
map beam, steering & F-M-Hp & speed quantised to 5 m/s & test & 1{,}917 & 37 & 105 & 88 & 84 & -3.3 [-5.0, -2.0] \\
map beam, steering & F-M-Hp & speed quantised to 5 m/s & few-shot & 1{,}959 & 21 & 59 & 92 & 85 & -6.7 [-8.7, -3.9] \\
map beam, steering & F-M-Hp & speed at 1 Hz & test & 1{,}917 & 37 & 37 & 88 & 87 & -0.2 [-1.1, +0.5] \\
map beam, steering & F-M-Hp & speed at 1 Hz & few-shot & 1{,}959 & 21 & 23 & 92 & 91 & -1.2 [-1.8, -0.6] \\
map beam, speed only & F-M-Hp & speed quantised to 1 m/s & test & 1{,}917 & 165 & 169 & 54 & 54 & -0.1 [-0.3, +0.1] \\
map beam, speed only & F-M-Hp & speed quantised to 1 m/s & few-shot & 1{,}959 & 343 & 338 & 47 & 47 & +0.1 [-0.4, +0.7] \\
map beam, speed only & F-M-Hp & speed quantised to 2 m/s & test & 1{,}917 & 165 & 174 & 54 & 54 & -0.1 [-0.5, +0.3] \\
map beam, speed only & F-M-Hp & speed quantised to 2 m/s & few-shot & 1{,}959 & 343 & 336 & 47 & 47 & -0.2 [-0.6, +0.4] \\
map beam, speed only & F-M-Hp & speed quantised to 5 m/s & test & 1{,}917 & 165 & 247 & 54 & 53 & -0.6 [-1.5, +0.1] \\
map beam, speed only & F-M-Hp & speed quantised to 5 m/s & few-shot & 1{,}959 & 343 & 323 & 47 & 47 & +0.5 [-0.5, +1.1] \\
map beam, speed only & F-M-Hp & speed at 1 Hz & test & 1{,}917 & 165 & 169 & 54 & 54 & -0.1 [-0.3, +0.1] \\
map beam, speed only & F-M-Hp & speed at 1 Hz & few-shot & 1{,}959 & 343 & 343 & 47 & 47 & +0.1 [+0.0, +0.3] \\
\bottomrule\end{tabular}
\end{table*}

\begin{figure}[t]\centering
\includegraphics[width=0.9\columnwidth]{fig_robust_w05.pdf}
\caption{Fixed map attacks on degraded inputs: change in the five-minute route recovery rate against the clean signal on
the same windows (percentage points; 95\,\% driver-clustered CIs; filled: test, hollow: few-shot drivers); every level of
the grid fixed before the first run; the 0.5\,s delay is a timing probe.}
\end{figure}
\begin{table*}[t]\centering\scriptsize\setlength{\tabcolsep}{3pt}
\caption{Degraded inputs for dead reckoning (no map): clean and degraded medians and the paired median / mean differences (metres).}\label{tab:supp_robust_dr}
\begin{tabular}{lllrrrrr}\toprule
method & degradation & fold & med. clean & med. degraded & $n$ & paired median [CI] & paired mean [CI] \\\midrule
steering DR & speed quantised to 1 m/s & test & 921 & 920 & 1{,}917 & -0.1 [-0.3, 0.2] & 0 [-1, 1] \\
steering DR & speed quantised to 1 m/s & few-shot & 1{,}669 & 1{,}664 & 1{,}963 & -0.1 [-0.6, 0.1] & -0 [-1, 1] \\
steering DR & speed noise 0.5 m/s & test & 921 & 926 & 1{,}917 & 0.1 [-0.0, 0.2] & 1 [0, 2] \\
steering DR & speed noise 0.5 m/s & few-shot & 1{,}669 & 1{,}670 & 1{,}963 & 0.1 [0.0, 0.3] & 2 [1, 3] \\
steering DR & steering at 1 Hz & test & 921 & 934 & 1{,}917 & 0.4 [-1.3, 1.6] & 7 [-1, 14] \\
steering DR & steering at 1 Hz & few-shot & 1{,}669 & 1{,}683 & 1{,}963 & 0.2 [-1.0, 2.3] & 5 [-1, 15] \\
straight line & speed quantised to 1 m/s & test & 735 & 736 & 1{,}917 & 0.0 [-0.1, 0.2] & 0 [0, 1] \\
straight line & speed quantised to 1 m/s & few-shot & 825 & 824 & 1{,}963 & 0.1 [-0.1, 0.5] & 1 [-0, 2] \\
straight line & speed noise 0.5 m/s & test & 735 & 736 & 1{,}917 & 0.1 [0.0, 0.3] & 1 [1, 2] \\
straight line & speed noise 0.5 m/s & few-shot & 825 & 825 & 1{,}963 & 0.1 [0.1, 0.2] & 1 [1, 1] \\
straight line & steering at 1 Hz & test & 735 & 735 & 1{,}917 & 0.0 [0.0, 0.0] & 0 [0, 0] \\
straight line & steering at 1 Hz & few-shot & 825 & 825 & 1{,}963 & 0.0 [0.0, 0.0] & 0 [0, 0] \\
\bottomrule\end{tabular}
\end{table*}

\begin{table*}[t]\centering\scriptsize\setlength{\tabcolsep}{3pt}
\caption{Cost study on one fixed validation subset, one worker per method: inference seconds per window (median and 90th percentile; failed searches count with the time they ran), peak memory and the per-drive graph build.}\label{tab:supp_cost}
\begin{tabular}{lrrrrrrrrrr}\toprule
method & $n_5$ & fail & med s & p90 s & $n_{15}$ & fail & med s & p90 s & RSS GB & graph s \\\midrule
dewri & 76 & 0 & 0.04 & 0.09 & 16 & 1 & 0.33 & 0.83 & 1.8 & 2.7 \\
dewri\_hp & 76 & 1 & 0.03 & 0.06 & 16 & 2 & 0.23 & 0.60 & 1.8 & 2.7 \\
dr\_none & 76 & 0 & 0.02 & 0.02 & 16 & 0 & 0.06 & 0.06 & 0.8 & -- \\
dr\_steer & 76 & 0 & 0.03 & 0.03 & 16 & 0 & 0.09 & 0.09 & 0.8 & -- \\
elastic\_pathing & 76 & 0 & 0.01 & 0.13 & 16 & 0 & 0.15 & 0.54 & 1.8 & 2.6 \\
elastic\_pathing\_hp & 76 & 0 & 0.01 & 0.12 & 16 & 0 & 0.16 & 0.52 & 1.8 & 2.6 \\
learned\_gru & 76 & 0 & 0.22 & 0.23 & 16 & 0 & 0.66 & 0.67 & 0.8 & -- \\
macrotracking & 76 & 0 & 0.08 & 0.21 & 16 & 0 & 0.43 & 0.69 & 1.8 & 2.6 \\
map\_beam\_none & 76 & 0 & 0.12 & 0.18 & 16 & 0 & 0.43 & 0.50 & 1.8 & 2.6 \\
map\_beam\_steer & 76 & 0 & 0.24 & 0.51 & 16 & 0 & 0.93 & 1.33 & 1.8 & 2.6 \\
map\_beam\_steer\_b50 & 76 & 0 & 0.07 & 0.13 & 16 & 0 & 0.25 & 0.37 & 1.8 & 2.6 \\
map\_pf\_learned & 76 & 0 & 0.83 & 0.95 & 16 & 0 & 2.56 & 2.89 & 1.7 & 2.6 \\
map\_pf\_none & 76 & 0 & 0.89 & 1.38 & 16 & 0 & 4.03 & 4.97 & 1.7 & 2.6 \\
map\_pf\_steer & 76 & 0 & 0.57 & 0.66 & 16 & 0 & 1.77 & 2.05 & 1.7 & 2.6 \\
map\_pf\_steer\_n500 & 76 & 0 & 0.22 & 0.30 & 16 & 0 & 0.70 & 0.87 & 1.7 & 2.6 \\
opd2 & 76 & 0 & 0.08 & 0.23 & 16 & 0 & 0.19 & 0.28 & 2.1 & 2.6 \\
spy & 76 & 0 & 0.13 & 0.68 & 16 & 0 & 0.24 & 1.23 & 1.8 & 2.7 \\
spy\_area & 76 & 9 & 4.93 & 13.53 & 16 & 1 & 7.71 & 23.97 & 1.8 & 2.7 \\
spy\_hp & 76 & 0 & 0.08 & 0.32 & 16 & 0 & 0.22 & 0.62 & 1.8 & 2.7 \\
\bottomrule\end{tabular}
\end{table*}

\begin{table*}[t]\centering\scriptsize\setlength{\tabcolsep}{3pt}
\caption{Accuracy on the cost subset at the published and at smaller search budgets (validation windows).}\label{tab:supp_cost_acc}
\begin{tabular}{lllrrrr}\toprule
method / budget & protocol & window & $n$ & fail & median ADE m & route \% \\\midrule
dewri\_\_G1 & F-M-Hm & 5 min & 76 & 0 & 1{,}233 & 21 \\
dewri\_\_G1 & F-M-Hm & 15 min & 16 & 1 & 4{,}170 & 12 \\
dewri\_hp\_\_G1 & F-M-Hp & 5 min & 76 & 1 & 524 & 34 \\
dewri\_hp\_\_G1 & F-M-Hp & 15 min & 16 & 2 & 1{,}737 & 25 \\
dr\_none\_\_G1 & O-N-Hp & 5 min & 76 & 0 & 497 & 16 \\
dr\_none\_\_G1 & O-N-Hp & 15 min & 16 & 0 & 3{,}132 & 6 \\
dr\_steer\_\_G3 & O-N-Hp & 5 min & 76 & 0 & 745 & 3 \\
dr\_steer\_\_G3 & O-N-Hp & 15 min & 16 & 0 & 4{,}681 & 0 \\
elastic\_pathing\_\_G1 & F-M-Hm & 5 min & 76 & 0 & 428 & 38 \\
elastic\_pathing\_\_G1 & F-M-Hm & 15 min & 16 & 0 & 3{,}268 & 19 \\
elastic\_pathing\_hp\_\_G1 & F-M-Hp & 5 min & 76 & 0 & 205 & 46 \\
elastic\_pathing\_hp\_\_G1 & F-M-Hp & 15 min & 16 & 0 & 2{,}100 & 25 \\
learned\_gru\_s0\_\_G3 & O-N-Hp & 5 min & 76 & 0 & 203 & 12 \\
learned\_gru\_s0\_\_G3 & O-N-Hp & 15 min & 16 & 0 & 1{,}237 & 0 \\
macrotracking\_\_G3 & O-M-Hp & 5 min & 76 & 0 & 75 & 50 \\
macrotracking\_\_G3 & O-M-Hp & 15 min & 16 & 0 & 1{,}709 & 25 \\
map\_beam\_none\_\_G1 & F-M-Hp & 5 min & 76 & 0 & 549 & 32 \\
map\_beam\_none\_\_G1 & F-M-Hp & 15 min & 16 & 0 & 4{,}260 & 12 \\
map\_beam\_steer\_\_G3 & F-M-Hp & 5 min & 76 & 0 & 21 & 84 \\
map\_beam\_steer\_\_G3 & F-M-Hp & 15 min & 16 & 0 & 58 & 81 \\
map\_beam\_steer\_b50\_\_G3 & F-M-Hp & 5 min & 76 & 0 & 21 & 84 \\
map\_beam\_steer\_b50\_\_G3 & F-M-Hp & 15 min & 16 & 0 & 58 & 81 \\
map\_pf\_learned\_gru\_s0\_\_G3 & O-M-Hp & 5 min & 76 & 0 & 38 & 84 \\
map\_pf\_learned\_gru\_s0\_\_G3 & O-M-Hp & 15 min & 16 & 0 & 115 & 88 \\
map\_pf\_none\_\_G1 & O-M-Hp & 5 min & 76 & 0 & 933 & 36 \\
map\_pf\_none\_\_G1 & O-M-Hp & 15 min & 16 & 0 & 3{,}935 & 19 \\
map\_pf\_steer\_\_G3 & O-M-Hp & 5 min & 76 & 0 & 26 & 87 \\
map\_pf\_steer\_\_G3 & O-M-Hp & 15 min & 16 & 0 & 90 & 81 \\
map\_pf\_steer\_n500\_\_G3 & O-M-Hp & 5 min & 76 & 0 & 24 & 87 \\
map\_pf\_steer\_n500\_\_G3 & O-M-Hp & 15 min & 16 & 0 & 91 & 81 \\
opd2\_\_G3 & F-M-Hp & 5 min & 76 & 0 & 23 & 76 \\
opd2\_\_G3 & F-M-Hp & 15 min & 16 & 0 & 1{,}317 & 31 \\
spy\_\_G3 & F-M-Hm & 5 min & 76 & 0 & 1{,}306 & 7 \\
spy\_\_G3 & F-M-Hm & 15 min & 16 & 0 & 5{,}401 & 6 \\
spy\_area\_\_G3 & F-M-Hm & 5 min & 76 & 9 & 2{,}743 & 0 \\
spy\_area\_\_G3 & F-M-Hm & 15 min & 16 & 1 & 6{,}180 & 0 \\
spy\_hp\_\_G3 & F-M-Hp & 5 min & 76 & 0 & 965 & 16 \\
spy\_hp\_\_G3 & F-M-Hp & 15 min & 16 & 0 & 4{,}946 & 6 \\
\bottomrule\end{tabular}
\end{table*}

\begin{table*}[t]\centering\scriptsize\setlength{\tabcolsep}{3pt}
\caption{Candidate confidence score of the offline beam: AUROC of the Top-1 score for route recovery (driver-clustered CI), Spearman correlation with the ADE, and the share of windows whose score is 1.}\label{tab:supp_conf}
\begin{tabular}{llllrrrrr}\toprule
method & protocol & window & fold & $n$ & fail & AUROC [CI] & Spearman & score = 1 \% \\\midrule
map\_beam\_steer & F-M-Hp & 5 min & test & 1{,}917 & 0 & 0.47 [0.41, 0.52] & -0.01 & 29 \\
map\_beam\_steer & F-M-Hp & 5 min & few-shot & 1{,}959 & 0 & 0.49 [0.44, 0.54] & -0.06 & 38 \\
map\_beam\_steer & F-M-Hp & 15 min & test & 608 & 0 & 0.54 [0.48, 0.59] & -0.08 & 19 \\
map\_beam\_steer & F-M-Hp & 15 min & few-shot & 561 & 0 & 0.54 [0.47, 0.59] & -0.20 & 34 \\
map\_beam\_steer & F-M-Hm & 5 min & test & 1{,}917 & 0 & 0.57 [0.53, 0.62] & -0.11 & 16 \\
map\_beam\_steer & F-M-Hm & 5 min & few-shot & 1{,}959 & 0 & 0.57 [0.51, 0.61] & -0.15 & 27 \\
map\_beam\_steer & F-M-Hm & 15 min & test & 608 & 0 & 0.55 [0.50, 0.60] & -0.09 & 15 \\
map\_beam\_steer & F-M-Hm & 15 min & few-shot & 561 & 0 & 0.58 [0.48, 0.65] & -0.26 & 30 \\
map\_beam\_none & F-M-Hp & 5 min & test & 1{,}917 & 0 & 0.81 [0.73, 0.86] & -0.51 & 7 \\
map\_beam\_none & F-M-Hp & 5 min & few-shot & 1{,}959 & 0 & 0.79 [0.74, 0.84] & -0.40 & 8 \\
map\_beam\_none & F-M-Hp & 15 min & test & 608 & 0 & 0.79 [0.63, 0.85] & -0.42 & 3 \\
map\_beam\_none & F-M-Hp & 15 min & few-shot & 561 & 0 & 0.77 [0.70, 0.82] & -0.26 & 2 \\
map\_beam\_none & F-M-Hm & 5 min & test & 1{,}917 & 0 & 0.65 [0.59, 0.67] & 0.02 & 1 \\
map\_beam\_none & F-M-Hm & 5 min & few-shot & 1{,}959 & 0 & 0.68 [0.65, 0.71] & 0.03 & 0 \\
map\_beam\_none & F-M-Hm & 15 min & test & 608 & 0 & 0.69 [0.56, 0.71] & 0.04 & 0 \\
map\_beam\_none & F-M-Hm & 15 min & few-shot & 561 & 0 & 0.69 [0.62, 0.73] & 0.07 & 0 \\
\bottomrule\end{tabular}
\end{table*}

\begin{table*}[t]\centering\scriptsize\setlength{\tabcolsep}{3pt}
\caption{Selective prediction by the beam's confidence score: windows accepted at a threshold (coverage, all expected windows in the denominator), route recovery among the accepted and the share of accepted routes that are wrong.}\label{tab:supp_conf_thr}
\begin{tabular}{llllrrrr}\toprule
method & protocol & window & fold & $\tau$ & coverage \% & route among accepted \% [CI] & wrong \% \\\midrule
map\_beam\_steer & F-M-Hp & 5 min & test & 0.5 & 77 & 86 [80, 92] & 14 \\
map\_beam\_steer & F-M-Hp & 5 min & test & 0.7 & 57 & 86 [79, 92] & 14 \\
map\_beam\_steer & F-M-Hp & 5 min & test & 0.9 & 46 & 87 [81, 92] & 13 \\
map\_beam\_steer & F-M-Hp & 5 min & few-shot & 0.5 & 85 & 92 [88, 94] & 8 \\
map\_beam\_steer & F-M-Hp & 5 min & few-shot & 0.7 & 65 & 91 [86, 94] & 9 \\
map\_beam\_steer & F-M-Hp & 5 min & few-shot & 0.9 & 55 & 92 [87, 94] & 8 \\
map\_beam\_steer & F-M-Hp & 15 min & test & 0.5 & 84 & 73 [59, 83] & 27 \\
map\_beam\_steer & F-M-Hp & 15 min & test & 0.7 & 61 & 74 [59, 85] & 26 \\
map\_beam\_steer & F-M-Hp & 15 min & test & 0.9 & 45 & 76 [61, 85] & 24 \\
map\_beam\_steer & F-M-Hp & 15 min & few-shot & 0.5 & 89 & 76 [70, 81] & 24 \\
map\_beam\_steer & F-M-Hp & 15 min & few-shot & 0.7 & 69 & 76 [68, 80] & 24 \\
map\_beam\_steer & F-M-Hp & 15 min & few-shot & 0.9 & 54 & 76 [70, 81] & 24 \\
map\_beam\_steer & F-M-Hm & 5 min & test & 0.5 & 70 & 83 [77, 88] & 17 \\
map\_beam\_steer & F-M-Hm & 5 min & test & 0.7 & 50 & 84 [78, 89] & 16 \\
map\_beam\_steer & F-M-Hm & 5 min & test & 0.9 & 36 & 85 [79, 91] & 15 \\
map\_beam\_steer & F-M-Hm & 5 min & few-shot & 0.5 & 80 & 87 [84, 90] & 13 \\
map\_beam\_steer & F-M-Hm & 5 min & few-shot & 0.7 & 61 & 88 [84, 90] & 12 \\
map\_beam\_steer & F-M-Hm & 5 min & few-shot & 0.9 & 48 & 89 [85, 92] & 11 \\
map\_beam\_steer & F-M-Hm & 15 min & test & 0.5 & 82 & 72 [58, 83] & 28 \\
map\_beam\_steer & F-M-Hm & 15 min & test & 0.7 & 58 & 73 [58, 84] & 27 \\
map\_beam\_steer & F-M-Hm & 15 min & test & 0.9 & 42 & 74 [58, 86] & 26 \\
map\_beam\_steer & F-M-Hm & 15 min & few-shot & 0.5 & 88 & 75 [70, 81] & 25 \\
map\_beam\_steer & F-M-Hm & 15 min & few-shot & 0.7 & 69 & 75 [68, 79] & 25 \\
map\_beam\_steer & F-M-Hm & 15 min & few-shot & 0.9 & 53 & 79 [72, 83] & 21 \\
map\_beam\_none & F-M-Hp & 5 min & test & 0.5 & 28 & 89 [72, 95] & 11 \\
map\_beam\_none & F-M-Hp & 5 min & test & 0.7 & 19 & 94 [84, 97] & 6 \\
map\_beam\_none & F-M-Hp & 5 min & test & 0.9 & 13 & 97 [89, 99] & 3 \\
map\_beam\_none & F-M-Hp & 5 min & few-shot & 0.5 & 25 & 82 [73, 90] & 18 \\
map\_beam\_none & F-M-Hp & 5 min & few-shot & 0.7 & 15 & 91 [83, 96] & 9 \\
map\_beam\_none & F-M-Hp & 5 min & few-shot & 0.9 & 11 & 94 [88, 98] & 6 \\
map\_beam\_none & F-M-Hp & 15 min & test & 0.5 & 20 & 86 [51, 95] & 14 \\
map\_beam\_none & F-M-Hp & 15 min & test & 0.7 & 13 & 91 [64, 96] & 9 \\
map\_beam\_none & F-M-Hp & 15 min & test & 0.9 & 7 & 93 [60, 99] & 7 \\
map\_beam\_none & F-M-Hp & 15 min & few-shot & 0.5 & 16 & 67 [42, 81] & 33 \\
map\_beam\_none & F-M-Hp & 15 min & few-shot & 0.7 & 8 & 84 [55, 94] & 16 \\
map\_beam\_none & F-M-Hp & 15 min & few-shot & 0.9 & 5 & 93 [76, 100] & 7 \\
map\_beam\_none & F-M-Hm & 5 min & test & 0.5 & 17 & 50 [39, 55] & 50 \\
map\_beam\_none & F-M-Hm & 5 min & test & 0.7 & 6 & 56 [37, 69] & 44 \\
map\_beam\_none & F-M-Hm & 5 min & test & 0.9 & 2 & 75 [40, 90] & 25 \\
map\_beam\_none & F-M-Hm & 5 min & few-shot & 0.5 & 14 & 46 [42, 53] & 54 \\
map\_beam\_none & F-M-Hm & 5 min & few-shot & 0.7 & 5 & 58 [51, 70] & 42 \\
map\_beam\_none & F-M-Hm & 5 min & few-shot & 0.9 & 2 & 62 [49, 78] & 38 \\
map\_beam\_none & F-M-Hm & 15 min & test & 0.5 & 18 & 50 [23, 58] & 50 \\
map\_beam\_none & F-M-Hm & 15 min & test & 0.7 & 9 & 58 [34, 71] & 42 \\
map\_beam\_none & F-M-Hm & 15 min & test & 0.9 & 2 & 57 [8, 94] & 43 \\
map\_beam\_none & F-M-Hm & 15 min & few-shot & 0.5 & 10 & 42 [26, 54] & 58 \\
map\_beam\_none & F-M-Hm & 15 min & few-shot & 0.7 & 5 & 50 [33, 70] & 50 \\
map\_beam\_none & F-M-Hm & 15 min & few-shot & 0.9 & 2 & 36 [12, 57] & 64 \\
\bottomrule\end{tabular}
\end{table*}

\begin{table*}[t]\centering\scriptsize\setlength{\tabcolsep}{3pt}
\caption{Window taxonomy (five minutes, whole-path criterion at the given threshold): no output, no candidate within the threshold, a lower-ranked candidate within it while the Top-1 is not (ranking failure), Top-1 within it.}\label{tab:supp_tax}
\begin{tabular}{lllrrrrrr}\toprule
method & protocol & fold & thr. m & $n$ & no output \% & no cand. \% & ranking \% & Top-1 ok \% \\\midrule
map\_beam\_steer & F-M-Hp & test & 50 & 1{,}917 & 0 & 28 & 10 & 61 \\
map\_beam\_steer & F-M-Hp & test & 100 & 1{,}917 & 0 & 12 & 7 & 81 \\
map\_beam\_steer & F-M-Hp & few-shot & 50 & 1{,}959 & 0 & 19 & 7 & 74 \\
map\_beam\_steer & F-M-Hp & few-shot & 100 & 1{,}959 & 0 & 6 & 6 & 88 \\
map\_beam\_steer & F-M-Hm & test & 50 & 1{,}917 & 0 & 29 & 15 & 57 \\
map\_beam\_steer & F-M-Hm & test & 100 & 1{,}917 & 0 & 13 & 12 & 75 \\
map\_beam\_steer & F-M-Hm & few-shot & 50 & 1{,}959 & 0 & 19 & 11 & 70 \\
map\_beam\_steer & F-M-Hm & few-shot & 100 & 1{,}959 & 0 & 7 & 10 & 83 \\
map\_beam\_none & F-M-Hp & test & 50 & 1{,}917 & 0 & 55 & 13 & 32 \\
map\_beam\_none & F-M-Hp & test & 100 & 1{,}917 & 0 & 41 & 13 & 46 \\
map\_beam\_none & F-M-Hp & few-shot & 50 & 1{,}959 & 0 & 53 & 18 & 29 \\
map\_beam\_none & F-M-Hp & few-shot & 100 & 1{,}959 & 0 & 40 & 18 & 42 \\
map\_beam\_none & F-M-Hm & test & 50 & 1{,}917 & 0 & 58 & 24 & 18 \\
map\_beam\_none & F-M-Hm & test & 100 & 1{,}917 & 0 & 45 & 28 & 27 \\
map\_beam\_none & F-M-Hm & few-shot & 50 & 1{,}959 & 0 & 57 & 27 & 16 \\
map\_beam\_none & F-M-Hm & few-shot & 100 & 1{,}959 & 0 & 45 & 30 & 25 \\
dewri\_hp & F-M-Hp & test & 50 & 1{,}917 & 1 & 45 & 18 & 37 \\
dewri\_hp & F-M-Hp & test & 100 & 1{,}917 & 1 & 25 & 20 & 54 \\
dewri\_hp & F-M-Hp & few-shot & 50 & 1{,}959 & 3 & 37 & 22 & 39 \\
dewri\_hp & F-M-Hp & few-shot & 100 & 1{,}959 & 3 & 22 & 23 & 52 \\
dewri & F-M-Hm & test & 50 & 1{,}917 & 0 & 49 & 22 & 29 \\
dewri & F-M-Hm & test & 100 & 1{,}917 & 0 & 30 & 28 & 41 \\
dewri & F-M-Hm & few-shot & 50 & 1{,}959 & 1 & 41 & 29 & 29 \\
dewri & F-M-Hm & few-shot & 100 & 1{,}959 & 1 & 27 & 32 & 39 \\
elastic\_pathing\_hp & F-M-Hp & test & 50 & 1{,}917 & 5 & 56 & 4 & 35 \\
elastic\_pathing\_hp & F-M-Hp & test & 100 & 1{,}917 & 5 & 40 & 3 & 52 \\
elastic\_pathing\_hp & F-M-Hp & few-shot & 50 & 1{,}959 & 4 & 52 & 5 & 39 \\
elastic\_pathing\_hp & F-M-Hp & few-shot & 100 & 1{,}959 & 4 & 37 & 5 & 54 \\
elastic\_pathing & F-M-Hm & test & 50 & 1{,}917 & 3 & 61 & 3 & 33 \\
elastic\_pathing & F-M-Hm & test & 100 & 1{,}917 & 3 & 46 & 2 & 48 \\
elastic\_pathing & F-M-Hm & few-shot & 50 & 1{,}959 & 2 & 62 & 4 & 32 \\
elastic\_pathing & F-M-Hm & few-shot & 100 & 1{,}959 & 2 & 49 & 4 & 45 \\
\bottomrule\end{tabular}
\end{table*}

\begin{table*}[t]\centering\scriptsize\setlength{\tabcolsep}{3pt}
\caption{Complementarity of the map methods (five-minute windows): per-window oracles over a method set (lowest ADE; any route recovered), the share of windows every member fails, and the fixed reference method.}\label{tab:supp_comp}
\begin{tabular}{llrrrrlrr}\toprule
set & fold & $n$ & oracle median m [CI] & oracle route \% [CI] & all fail \% [CI] & reference & ref. median & ref. route \% \\\midrule
A steering map, heading known & test & 1{,}917 & 25 [18, 35] & 97 [95, 98] & 3.1 [2.1, 5.2] & beam & 37 & 88 \\
A steering map, heading known & few-shot & 1{,}959 & 15 [11, 28] & 98 [97, 99] & 1.9 [1.2, 2.8] & beam & 21 & 92 \\
B speed-only map, heading known & test & 1{,}917 & 50 [25, 93] & 78 [68, 86] & 22.3 [14.3, 32.5] & beam\_none & 165 & 54 \\
B speed-only map, heading known & few-shot & 1{,}959 & 40 [21, 67] & 77 [70, 82] & 22.8 [18.2, 29.9] & beam\_none & 343 & 47 \\
B' speed-only map, heading unknown & test & 1{,}917 & 61 [29, 139] & 72 [59, 81] & 28.4 [19.0, 40.5] & beam\_none\_hm & 1{,}314 & 32 \\
B' speed-only map, heading unknown & few-shot & 1{,}959 & 54 [36, 81] & 70 [64, 74] & 30.1 [25.7, 36.2] & beam\_none\_hm & 1{,}527 & 29 \\
C published, heading known & test & 1{,}917 & 30 [21, 46] & 94 [90, 95] & 6.3 [4.9, 9.5] & nan & -- & -- \\
C published, heading known & few-shot & 1{,}959 & 17 [13, 36] & 95 [93, 97] & 4.7 [3.4, 7.5] & nan & -- & -- \\
D all map, heading known & test & 1{,}917 & 24 [18, 34] & 98 [96, 99] & 2.3 [1.4, 3.6] & beam & 37 & 88 \\
D all map, heading known & few-shot & 1{,}959 & 14 [11, 26] & 98 [97, 99] & 1.6 [0.9, 2.5] & beam & 21 & 92 \\
P pedal subset, heading unknown & test & 245 & 40 [29, 70] & 95 [86, 97] & 4.9 [2.7, 14.1] & beam\_hm & 46 & 88 \\
P pedal subset, heading unknown & few-shot & 162 & 19 [14, 35] & 93 [87, 98] & 7.4 [2.4, 13.3] & beam\_hm & 21 & 86 \\
\bottomrule\end{tabular}
\end{table*}

\begin{table*}[t]\centering\scriptsize\setlength{\tabcolsep}{3pt}
\caption{Common failure of the steering and speed map methods by reference turn count and road density (folds pooled).}\label{tab:supp_cfail}
\begin{tabular}{llrrrrr}\toprule
stratum & level & $n$ & drivers & all fail \% [CI] & oracle route \% & beam route \% \\\midrule
reference turns & 0-1 & 2{,}234 & 81 & 0.1 [0.0, 0.2] & 100 & 98 \\
reference turns & 2-3 & 712 & 72 & 1.4 [0.6, 2.9] & 99 & 88 \\
reference turns & 4-6 & 528 & 74 & 5.3 [3.3, 8.7] & 95 & 73 \\
reference turns & 7+ & 402 & 66 & 8.7 [5.3, 12.4] & 91 & 67 \\
density & low (56/246 per km2 cuts) & 1{,}292 & 33 & 0.9 [0.2, 1.7] & 99 & 97 \\
density & mid (56/246 per km2 cuts) & 1{,}292 & 82 & 1.9 [1.2, 2.8] & 98 & 87 \\
density & high (56/246 per km2 cuts) & 1{,}292 & 81 & 2.9 [1.6, 4.5] & 97 & 86 \\
all & all & 3{,}876 & 92 & 1.9 [1.4, 2.6] & 98 & 90 \\
\bottomrule\end{tabular}
\end{table*}

\begin{table*}[t]\centering\scriptsize\setlength{\tabcolsep}{3pt}
\caption{Per-driver route recovery share (five-minute windows; drivers with at least five windows): pooled rate, quantiles over drivers and the share of drivers below 50\,\%.}\label{tab:supp_driver}
\begin{tabular}{llrrrrrrrrrr}\toprule
method & fold & windows & drivers & $\ge 5$ & pooled \% & q10 & q25 & q50 & q75 & q90 & below 50 \% \\\midrule
pf & test & 1{,}917 & 42 & 36 & 87 & 62 & 77 & 87 & 96 & 100 & 3 \\
pf & few-shot & 1{,}959 & 50 & 41 & 90 & 73 & 78 & 92 & 96 & 100 & 0 \\
beam & test & 1{,}917 & 42 & 36 & 88 & 62 & 75 & 88 & 100 & 100 & 3 \\
beam & few-shot & 1{,}959 & 50 & 41 & 92 & 77 & 83 & 94 & 97 & 100 & 0 \\
opd2 & test & 1{,}917 & 42 & 36 & 79 & 54 & 64 & 78 & 84 & 89 & 8 \\
opd2 & few-shot & 1{,}959 & 50 & 41 & 78 & 54 & 67 & 79 & 88 & 91 & 5 \\
macro & test & 1{,}917 & 42 & 36 & 62 & 21 & 33 & 50 & 73 & 82 & 50 \\
macro & few-shot & 1{,}959 & 50 & 41 & 65 & 32 & 42 & 61 & 77 & 84 & 37 \\
ep\_hp & test & 1{,}917 & 42 & 36 & 58 & 23 & 31 & 41 & 64 & 71 & 61 \\
ep\_hp & few-shot & 1{,}959 & 50 & 41 & 59 & 30 & 36 & 50 & 67 & 71 & 46 \\
dewri\_hp & test & 1{,}917 & 42 & 36 & 62 & 15 & 26 & 45 & 61 & 78 & 53 \\
dewri\_hp & few-shot & 1{,}959 & 50 & 41 & 59 & 17 & 30 & 56 & 70 & 85 & 49 \\
spy\_hp & test & 1{,}917 & 42 & 36 & 20 & 0 & 0 & 10 & 18 & 33 & 100 \\
spy\_hp & few-shot & 1{,}959 & 50 & 41 & 22 & 0 & 4 & 14 & 27 & 35 & 100 \\
pf\_none & test & 1{,}917 & 42 & 36 & 56 & 9 & 22 & 35 & 58 & 66 & 61 \\
pf\_none & few-shot & 1{,}959 & 50 & 41 & 48 & 12 & 29 & 50 & 64 & 75 & 49 \\
beam\_none & test & 1{,}917 & 42 & 36 & 54 & 8 & 20 & 31 & 57 & 67 & 64 \\
beam\_none & few-shot & 1{,}959 & 50 & 41 & 47 & 15 & 30 & 50 & 62 & 71 & 49 \\
ep & test & 1{,}917 & 42 & 36 & 53 & 14 & 25 & 34 & 51 & 69 & 72 \\
ep & few-shot & 1{,}959 & 50 & 41 & 50 & 23 & 32 & 43 & 57 & 67 & 59 \\
dewri & test & 1{,}917 & 42 & 36 & 46 & 9 & 17 & 35 & 50 & 63 & 72 \\
dewri & few-shot & 1{,}959 & 50 & 41 & 43 & 17 & 21 & 38 & 47 & 63 & 78 \\
\bottomrule\end{tabular}
\end{table*}

\begin{table*}[t]\centering\scriptsize\setlength{\tabcolsep}{3pt}
\caption{Road-density strata around the start (tertiles; folds pooled): median ADE with CI, 100\,m coverage and tail-lost share.}\label{tab:supp_density}
\begin{tabular}{lllrrrrrr}\toprule
method & protocol & window & tertile & $n$ & drivers & median ADE m [CI] & c100 \% & tail lost \% \\\midrule
map\_pf\_steer\_\_G3 & O-M-Hp & 5 min & low (< 56/km2) & 1{,}292 & 33 & 29 [18, 59] & 90 & 17 \\
map\_pf\_steer\_\_G3 & O-M-Hp & 5 min & mid (56 to 246/km2) & 1{,}292 & 82 & 39 [27, 48] & 83 & 27 \\
map\_pf\_steer\_\_G3 & O-M-Hp & 5 min & high (> 246/km2) & 1{,}292 & 81 & 36 [26, 45] & 83 & 27 \\
map\_pf\_steer\_\_G3 & O-M-Hp & 15 min & low (< 56/km2) & 414 & 27 & 89 [46, 185] & 65 & 53 \\
map\_pf\_steer\_\_G3 & O-M-Hp & 15 min & mid (56 to 246/km2) & 381 & 65 & 202 [132, 373] & 50 & 70 \\
map\_pf\_steer\_\_G3 & O-M-Hp & 15 min & high (> 246/km2) & 374 & 61 & 264 [169, 564] & 51 & 67 \\
map\_beam\_steer\_\_G3 & F-M-Hp & 5 min & low (< 56/km2) & 1{,}292 & 33 & 24 [16, 53] & 91 & 17 \\
map\_beam\_steer\_\_G3 & F-M-Hp & 5 min & mid (56 to 246/km2) & 1{,}292 & 82 & 32 [23, 42] & 85 & 25 \\
map\_beam\_steer\_\_G3 & F-M-Hp & 5 min & high (> 246/km2) & 1{,}292 & 81 & 33 [22, 40] & 84 & 27 \\
map\_beam\_steer\_\_G3 & F-M-Hp & 15 min & low (< 56/km2) & 414 & 27 & 77 [44, 167] & 67 & 53 \\
map\_beam\_steer\_\_G3 & F-M-Hp & 15 min & mid (56 to 246/km2) & 381 & 65 & 163 [105, 326] & 53 & 70 \\
map\_beam\_steer\_\_G3 & F-M-Hp & 15 min & high (> 246/km2) & 374 & 61 & 162 [100, 275] & 57 & 61 \\
map\_pf\_none\_\_G1 & O-M-Hp & 5 min & low (< 56/km2) & 1{,}292 & 33 & 73 [32, 537] & 73 & 37 \\
map\_pf\_none\_\_G1 & O-M-Hp & 5 min & mid (56 to 246/km2) & 1{,}292 & 82 & 677 [234, 953] & 49 & 62 \\
map\_pf\_none\_\_G1 & O-M-Hp & 5 min & high (> 246/km2) & 1{,}292 & 81 & 843 [627, 978] & 40 & 70 \\
map\_pf\_none\_\_G1 & O-M-Hp & 15 min & low (< 56/km2) & 414 & 27 & 1{,}023 [96, 5{,}289] & 42 & 77 \\
map\_pf\_none\_\_G1 & O-M-Hp & 15 min & mid (56 to 246/km2) & 381 & 65 & 3{,}457 [2{,}459, 4{,}268] & 23 & 86 \\
map\_pf\_none\_\_G1 & O-M-Hp & 15 min & high (> 246/km2) & 374 & 61 & 3{,}468 [2{,}920, 4{,}276] & 19 & 86 \\
dr\_steer\_\_G3 & O-N-Hp & 5 min & low (< 56/km2) & 1{,}296 & 33 & 2{,}102 [1{,}067, 4{,}256] & 15 & 99 \\
dr\_steer\_\_G3 & O-N-Hp & 5 min & mid (56 to 246/km2) & 1{,}292 & 82 & 1{,}245 [819, 1{,}859] & 22 & 95 \\
dr\_steer\_\_G3 & O-N-Hp & 5 min & high (> 246/km2) & 1{,}292 & 81 & 632 [428, 1{,}132] & 30 & 90 \\
dr\_steer\_\_G3 & O-N-Hp & 15 min & low (< 56/km2) & 418 & 27 & 8{,}506 [6{,}751, 12{,}469] & 5 & 99 \\
dr\_steer\_\_G3 & O-N-Hp & 15 min & mid (56 to 246/km2) & 381 & 65 & 6{,}461 [4{,}875, 7{,}873] & 9 & 99 \\
dr\_steer\_\_G3 & O-N-Hp & 15 min & high (> 246/km2) & 374 & 61 & 3{,}373 [1{,}772, 5{,}018] & 13 & 97 \\
dr\_curvature\_\_G5 & O-N-Hp & 5 min & low (< 56/km2) & 1{,}296 & 33 & 85 [74, 109] & 65 & 67 \\
dr\_curvature\_\_G5 & O-N-Hp & 5 min & mid (56 to 246/km2) & 1{,}292 & 82 & 86 [76, 102] & 63 & 63 \\
dr\_curvature\_\_G5 & O-N-Hp & 5 min & high (> 246/km2) & 1{,}292 & 81 & 70 [64, 82] & 69 & 55 \\
dr\_curvature\_\_G5 & O-N-Hp & 15 min & low (< 56/km2) & 418 & 27 & 516 [479, 724] & 25 & 94 \\
dr\_curvature\_\_G5 & O-N-Hp & 15 min & mid (56 to 246/km2) & 381 & 65 & 548 [443, 675] & 25 & 95 \\
dr\_curvature\_\_G5 & O-N-Hp & 15 min & high (> 246/km2) & 374 & 61 & 418 [350, 492] & 31 & 89 \\
\bottomrule\end{tabular}
\end{table*}

\begin{table*}[t]\centering\scriptsize\setlength{\tabcolsep}{3pt}
\caption{External check on the controlled supplement (one vehicle, one metropolitan area, other drivers; five-minute windows), methods unchanged after freezing.}\label{tab:supp_controlled}
\begin{tabular}{llrrrrrrr}\toprule
method & protocol & $n$ & drives & drivers & fail & median ADE m [CI] & p90 m & route \% [CI] \\\midrule
steering DR, no map & O-N-Hp & 46 & 27 & 12 & 0 & 117 [93, 170] & 449 & 7 [0, 10] \\
map PF, steering & O-M-Hp & 46 & 27 & 12 & 0 & 75 [46, 492] & 896 & 59 [26, 69] \\
MacroTracking & O-M-Hp & 46 & 27 & 12 & 0 & 19 [14, 24] & 144 & 85 [74, 95] \\
map beam, steering & F-M-Hp & 46 & 27 & 12 & 0 & 46 [34, 308] & 846 & 61 [37, 67] \\
OPD-II & F-M-Hp & 46 & 27 & 12 & 0 & 25 [11, 270] & 600 & 65 [32, 75] \\
map PF, speed & O-M-Hp & 46 & 27 & 12 & 0 & 762 [503, 1{,}011] & 1{,}545 & 9 [0, 20] \\
map beam, speed & F-M-Hp & 46 & 27 & 12 & 0 & 763 [321, 890] & 1{,}653 & 4 [0, 14] \\
Elastic Pathing, heading & F-M-Hp & 46 & 27 & 12 & 0 & 337 [239, 779] & 1{,}185 & 28 [6, 38] \\
Dewri et al., heading & F-M-Hp & 46 & 27 & 12 & 0 & 605 [345, 670] & 1{,}621 & 15 [0, 27] \\
\bottomrule\end{tabular}
\end{table*}

\begin{table*}[t]\centering\scriptsize\setlength{\tabcolsep}{3pt}
\caption{Map-extent checks: share of windows whose reference route leaves the 5\,km area adaptation square or the visible disc.}\label{tab:supp_area}
\begin{tabular}{llllrrrr}\toprule
check & method & window & fold & $n$ & out & share \% & median max m \\\midrule
out\_square & ct & 5 min & test & 245 & 181 & 73.9 & 4{,}486 \\
out\_square & ct & 5 min & few-shot & 162 & 106 & 65.4 & 3{,}822 \\
out\_square & spy & 5 min & test & 1{,}917 & 1{,}531 & 79.9 & 5{,}728 \\
out\_square & spy & 5 min & few-shot & 1{,}959 & 1{,}544 & 78.8 & 4{,}702 \\
out\_disc & G3 & 5 min & test & 1{,}917 & 0 & 0.0 & 6{,}311 \\
out\_disc & G3 & 5 min & few-shot & 1{,}959 & 0 & 0.0 & 5{,}168 \\
out\_disc & G3 & 15 min & test & 608 & 0 & 0.0 & 17{,}222 \\
out\_disc & G3 & 15 min & few-shot & 561 & 0 & 0.0 & 14{,}538 \\
out\_disc & G3 & whole drive & test & 282 & 18 & 6.4 & 10{,}190 \\
out\_disc & G3 & whole drive & few-shot & 454 & 37 & 8.1 & 6{,}263 \\
out\_disc & G1 & 5 min & test & 1{,}917 & 0 & 0.0 & 6{,}311 \\
out\_disc & G1 & 5 min & few-shot & 1{,}959 & 0 & 0.0 & 5{,}168 \\
out\_disc & G1 & 15 min & test & 608 & 0 & 0.0 & 17{,}222 \\
out\_disc & G1 & 15 min & few-shot & 561 & 0 & 0.0 & 14{,}538 \\
out\_disc & G1 & whole drive & test & 282 & 18 & 6.4 & 10{,}190 \\
out\_disc & G1 & whole drive & few-shot & 454 & 37 & 8.1 & 6{,}263 \\
\bottomrule\end{tabular}
\end{table*}


\section{S5. Agent study}
The supplementary agent study (Section VI of the paper) compares, on one five-minute window per held-out drive
(\AgentCohortTestN{} test windows of \AgentCohortTestDrivers{} drivers, \AgentCohortFewShotN{} few-shot windows of
\AgentCohortFewShotDrivers{} drivers), the offline beam (A), a supervised selector over the three reconstructors (B)
and a Qwen3-14B tool-calling controller (C). The protocol was frozen after a validation pilot
(\texttt{reports/agent\_protocol\_v1.md}): model revision, sampling, prompt, tool schema, budget, selector weights
and cohort hashes are recorded there; every window of every seed received a valid choice. Tables~\ref{tab:supp_agent}
to~\ref{tab:supp_agent_pilot} give every row.
\begin{table*}[t]\centering\scriptsize\setlength{\tabcolsep}{3pt}
\caption{Agent study (supplementary cohort, one five-minute window per evaluation drive): every configuration and seed, with the selector fallback row and the pool oracle.}\label{tab:supp_agent}
\begin{tabular}{p{0.26\textwidth}lrrrrrrr}\toprule
row & fold & $n$ & route \% [CI] & median ADE m [CI] & no out. \% & cost s & calls & tokens \\\midrule
map beam (A) & test & 268 & 69 [58, 85] & 49 [29, 66] & 0.0 & 0.3 & -- & -- \\
supervised selector (B) & test & 268 & 75 [67, 87] & 38 [29, 50] & 0.0 & 1.1 & 3.0 & -- \\
Qwen agent, seed 0 & test & 268 & 75 [66, 88] & 41 [29, 55] & 0.0 & 35.5 & 2.9 & 1{,}703 \\
Qwen agent, seed 1 & test & 268 & 75 [65, 87] & 44 [27, 54] & 0.0 & 37.8 & 2.9 & 1{,}753 \\
Qwen agent, seed 2 & test & 268 & 75 [66, 88] & 39 [29, 52] & 0.0 & 35.9 & 3.0 & 1{,}702 \\
Qwen agent (C), seeds pooled & test & 268 & 75 [66, 88] & 41 [29, 53] & 0.0 & 36.4 & 2.9 & 1{,}719 \\
Qwen agent with selector fallback & test & 268 & 75 [66, 88] & 41 [29, 53] & 0.0 & -- & -- & -- \\
pool oracle (diagnostic) & test & 268 & 83 [74, 93] & 27 [19, 32] & 0.0 & -- & -- & -- \\
map beam (A) & few-shot & 391 & 83 [76, 90] & 24 [18, 36] & 0.0 & 0.3 & -- & -- \\
supervised selector (B) & few-shot & 391 & 83 [78, 90] & 24 [18, 34] & 0.0 & 1.0 & 3.0 & -- \\
Qwen agent, seed 0 & few-shot & 391 & 84 [77, 91] & 23 [18, 34] & 0.0 & 31.9 & 3.0 & 1{,}626 \\
Qwen agent, seed 1 & few-shot & 391 & 84 [77, 90] & 23 [18, 34] & 0.0 & 37.9 & 2.9 & 1{,}741 \\
Qwen agent, seed 2 & few-shot & 391 & 84 [78, 91] & 22 [18, 35] & 0.0 & 35.2 & 2.9 & 1{,}711 \\
Qwen agent (C), seeds pooled & few-shot & 391 & 84 [77, 91] & 23 [18, 34] & 0.0 & 35.0 & 2.9 & 1{,}693 \\
Qwen agent with selector fallback & few-shot & 391 & 84 [77, 91] & 23 [18, 34] & 0.0 & -- & -- & -- \\
pool oracle (diagnostic) & few-shot & 391 & 92 [87, 96] & 16 [11, 23] & 0.0 & -- & -- & -- \\
\bottomrule\end{tabular}
\end{table*}

\begin{table*}[t]\centering\scriptsize\setlength{\tabcolsep}{3pt}
\caption{Agent study, paired comparisons on the same windows: route recovery difference in percentage points (seeds pooled inside the driver resamples; per-seed values) and the median ADE difference on windows with output on both sides.}\label{tab:supp_agent_pairs}
\begin{tabular}{llrrrr}\toprule
pair & fold & $n$ & $\Delta$ route pp [CI] & per seed & $\Delta$ median ADE m [CI] \\\midrule
agent\_vs\_selector & test & 268 & +0.1 [-1.6, +3.4] & +0.4, -0.4, +0.4 & 0.0 [0.0, 0.0] \\
agent\_vs\_beam & test & 268 & +5.7 [+1.6, +8.5] & +6.0, +5.2, +6.0 & 0.0 [0.0, 0.0] \\
agent\_vs\_selector & few-shot & 391 & +0.7 [-2.6, +2.8] & +0.8, +0.5, +0.8 & 0.0 [0.0, 0.0] \\
agent\_vs\_beam & few-shot & 391 & +1.2 [-1.3, +3.8] & +1.3, +1.0, +1.3 & 0.0 [0.0, 0.0] \\
\bottomrule\end{tabular}
\end{table*}

\begin{table}[H]\centering\scriptsize\setlength{\tabcolsep}{3pt}
\caption{Supervised selector on the validation cohort (cross-fitted by driver).}\label{tab:supp_agent_sel}
\begin{tabular}{lrrr}\toprule
row & windows & drivers & route \% [CI] \\\midrule
selector\_cv & 610 & 43 & 82 [77, 87] \\
beam1 & 610 & 43 & 79 [72, 84] \\
pf1 & 610 & 43 & 75 [69, 80] \\
ep1 & 610 & 43 & 28 [21, 40] \\
oracle & 610 & 43 & 91 [88, 94] \\
\bottomrule\end{tabular}
\end{table}

\begin{table*}[t]\centering\scriptsize\setlength{\tabcolsep}{3pt}
\caption{Validation pilot of the agent (40 windows, three seeds).}\label{tab:supp_agent_pilot}
\begin{tabular}{lrrrrrrrrr}\toprule
tag & $n$ & valid \% & trunc. \% & forced \% & turns & calls & tokens & thinking & LLM s \\\midrule
agent\_qwen3\_s0 & 40 & 100.0 & 2.5 & 87.5 & 3.8 & 3.0 & 1{,}703 & 1{,}521 & 27.9 \\
agent\_qwen3\_s1 & 40 & 100.0 & 0.0 & 85.0 & 3.8 & 2.9 & 1{,}754 & 1{,}577 & 38.0 \\
agent\_qwen3\_s2 & 40 & 100.0 & 0.0 & 77.5 & 3.7 & 2.9 & 1{,}693 & 1{,}513 & 35.6 \\
all & 120 & 100.0 & 0.8 & 83.3 & 3.8 & 2.9 & 1{,}717 & 1{,}537 & 34.6 \\
\bottomrule\end{tabular}
\end{table*}

\paragraph{Cost} The reconstructed serial cost of a window is the model's wall time plus the measured single-worker time
of the tools it ran plus the tool execution time. On 10 validation windows every reconstructor call was executed uncached and single-stream: median end-to-end time 27.7\,s against a reconstructed serial cost of 23.4\,s on the same windows (ratio 1.08, 90th percentile 1.68).


\paragraph{System prompt (verbatim)}
\begin{scriptsize}\begin{verbatim}
You are the controller of a vehicle-trajectory recovery
system. One five-minute driving window is given. The
system knows only the vehicle's CAN signals (speed and
steering-wheel angle), the start position (the origin),
the start heading, static vehicle parameters and the road
graph within the visible disc around the start. No GPS
positions, no true route and no scores are available to
you or to the tools; your job is to pick the candidate
path most likely to be the vehicle's true path.

Tools: run_reconstructor(name) runs one of three existing
reconstructors (beam, pf, ep) and returns its ranked
candidate paths with diagnostics;
inspect_candidate(candidate_id) gives the full diagnostics
of a candidate you have already seen;
query_local_graph(candidate_id, distance_m) describes the
road junction nearest to that distance along a candidate
you have seen; submit_choice(candidate_id, reason) ends
the task. Every reconstructor can be run at most once. You
must finish with submit_choice for a candidate id that a
reconstructor returned to you; nothing else is accepted.

Diagnostics (all computed from the CAN signals, the given
start heading and the road graph): path_len_m is the
candidate's length; len_ratio_dev is |path length /
CAN-integrated distance - 1|; hd_abs_rms_deg is the RMS
difference between the candidate's road bearing and the
dead-reckoned heading (start heading plus integrated
steering) over the moving seconds; hd_drift_rms_deg is the
same after removing the best constant offset; turns_path
and turns_dr count turns (bearing change above 45 degrees
within 40 m) along the candidate and in the dead-reckoned
heading, turns_unmatched counts turns present in one but
not the other (same sign, within 50 m of path distance);
uturns counts reversals; overlap_with_called_top1 is the
share of the candidate within 15 m of the Top-1 of another
reconstructor you have run (unknown for reconstructors you
have not run); end_dist_to_called_top1_m compares end
points; n_called_tools_agree counts other run
reconstructors whose Top-1 overlaps this candidate by at
least 0.8. Positions are metres east and north of the
start; compass headings are degrees clockwise from north.
Dead reckoning accumulates error, so small heading or turn
disagreement does not by itself rule out a candidate;
large unmatched turns or a path much shorter or longer
than the CAN distance are strong evidence against it.

Budget: at most 4 decision turns, at most 7 tool calls
before the final submit_choice, and a limited number of
generated tokens. You may request several tool calls in
one turn (for example two reconstructors at once); each
call counts toward the 7. Be concise. When the budget
forces it, you will be asked to submit immediately.

\end{verbatim}\end{scriptsize}

\paragraph{Tool schema and budget}
\begin{scriptsize}\begin{verbatim}
run_reconstructor(name)
Run one trajectory reconstructor on this window (each
reconstructor can be run once; a repeat returns the same
result and still counts). Returns its ranked candidate
paths with observable diagnostics.
inspect_candidate(candidate_id)
Detailed diagnostics of one candidate you have already
seen: heading consistency, turn matching against the
dead-reckoned heading, end position, overlap with the
Top-1 of the other reconstructors you have run, and where
it diverges from them.
query_local_graph(candidate_id, distance_m)
Road junction nearest to a given distance along a
candidate you have seen: arrival heading, the exits
(relative turn, length, road class, one-way), which exit
the candidate took, and the dead-reckoned heading change
around that moment.
submit_choice(candidate_id, reason)
Final answer: the id of one candidate you have seen. The
program outputs that candidate's full trajectory.
budget: 4 decision turns, 7 tool calls + 1 submit, 4096
completion tokens (per turn 2048, reserve 256), virtual
clock 600 s
\end{verbatim}\end{scriptsize}

\IfFileExists{tables/agent_max_rows.tex}{\paragraph{Extended configuration (B$^{+}$ and C$^{+}$ in Table IV of the paper)} After the frozen study, the same controller was given ten reconstructor
presets (wider and looser beams, four more filter seeds, Elastic Pathing without heading), two further tools (side-by-side
comparison with the divergence junction; the selector's probability), dead-reckoning shape and junction-consistency
evidence, calibration hints from the validation drives and a relaxed budget (12 turns, 40 calls, 24{,}576 tokens; the
served context was 16{,}384 tokens); a critic pass and a three-seed vote were added as variants. The configuration was
specified after the frozen study from validation data only and run once on the same 659 windows; the rows below also
pair it with the frozen configuration. These runs are exploratory and reported in full.
\begin{table*}[t]\centering\scriptsize\setlength{\tabcolsep}{3pt}
\caption{Ceiling exploration of the agent (C-max: ten reconstructor presets, comparison and scoring tools, relaxed budgets, optional critic and seed vote); exploratory, not part of the frozen protocol.}\label{tab:supp_agent_max}
\begin{tabular}{llrrrrrrr}\toprule
row & fold & $n$ & route \% [CI] & median ADE m [CI] & no out. \% & cost s & calls & tokens \\\midrule
map beam (A) & test & 268 & 69 [58, 85] & 49 [29, 66] & 0.0 & 0.3 & -- & -- \\
supervised selector (B) & test & 268 & 74 [66, 85] & 43 [28, 59] & 0.0 & 4.8 & 3.0 & -- \\
Qwen agent, seed 0 & test & 268 & 77 [70, 88] & 41 [28, 52] & 0.4 & 116.6 & 4.8 & 3{,}448 \\
Qwen agent, seed 1 & test & 268 & 79 [72, 89] & 38 [27, 49] & 0.4 & 116.8 & 4.9 & 3{,}497 \\
Qwen agent, seed 2 & test & 268 & 78 [70, 89] & 39 [27, 46] & 0.0 & 73.7 & 4.8 & 3{,}420 \\
Qwen agent (C), seeds pooled & test & 268 & 78 [71, 88] & 39 [27, 49] & 0.2 & 102.4 & 4.8 & 3{,}455 \\
Qwen agent with selector fallback & test & 268 & 78 [71, 88] & 39 [27, 49] & 0.0 & -- & -- & -- \\
pool oracle (diagnostic) & test & 268 & 87 [83, 94] & 24 [18, 29] & 0.0 & -- & -- & -- \\
agent\_qwen3max\_vote & test & 268 & 79 [72, 89] & 37 [27, 45] & 0.0 & -- & -- & -- \\
agent\_qwen3maxcritic\_s0 & test & 268 & 78 [71, 88] & 40 [27, 50] & 0.0 & -- & -- & -- \\
agent of run icde\_agent (seeds pooled) & test & 268 & 75 [66, 88] & 41 [29, 53] & 0.0 & -- & -- & -- \\
map beam (A) & few-shot & 391 & 83 [76, 90] & 24 [18, 36] & 0.0 & 0.3 & -- & -- \\
supervised selector (B) & few-shot & 391 & 82 [75, 90] & 26 [20, 41] & 0.0 & 4.4 & 3.0 & -- \\
Qwen agent, seed 0 & few-shot & 391 & 81 [75, 87] & 27 [20, 39] & 0.0 & 115.2 & 4.7 & 3{,}559 \\
Qwen agent, seed 1 & few-shot & 391 & 83 [76, 89] & 26 [20, 37] & 0.0 & 109.9 & 4.8 & 3{,}498 \\
Qwen agent, seed 2 & few-shot & 391 & 81 [75, 88] & 28 [21, 37] & 0.3 & 72.5 & 4.7 & 3{,}422 \\
Qwen agent (C), seeds pooled & few-shot & 391 & 81 [75, 88] & 27 [21, 37] & 0.1 & 99.2 & 4.8 & 3{,}493 \\
Qwen agent with selector fallback & few-shot & 391 & 81 [75, 88] & 27 [21, 37] & 0.0 & -- & -- & -- \\
pool oracle (diagnostic) & few-shot & 391 & 94 [90, 97] & 15 [11, 22] & 0.0 & -- & -- & -- \\
agent\_qwen3max\_vote & few-shot & 391 & 83 [76, 89] & 25 [20, 36] & 0.0 & -- & -- & -- \\
agent\_qwen3maxcritic\_s0 & few-shot & 391 & 82 [76, 88] & 26 [20, 34] & 0.3 & -- & -- & -- \\
agent of run icde\_agent (seeds pooled) & few-shot & 391 & 84 [77, 91] & 23 [18, 34] & 0.0 & -- & -- & -- \\
agent, both folds pooled & both folds & 659 & 80 [75, 85] & 31 [25, 39] & 0.2 & -- & -- & -- \\
agent\_prev, both folds pooled & both folds & 659 & 80 [74, 87] & 30 [22, 41] & 0.0 & -- & -- & -- \\
selector, both folds pooled & both folds & 659 & 79 [73, 85] & 31 [25, 44] & 0.0 & -- & -- & -- \\
beam, both folds pooled & both folds & 659 & 77 [69, 86] & 32 [23, 47] & 0.0 & -- & -- & -- \\
\bottomrule\end{tabular}
\end{table*}

\begin{table}[H]\centering\scriptsize\setlength{\tabcolsep}{3pt}
\caption{Ceiling exploration, paired route-recovery differences (percentage points, driver-clustered CIs).}\label{tab:supp_agent_max_pairs}
\begin{tabular}{llrr}\toprule
pair & fold & $n$ & $\Delta$ route pp [CI] \\\midrule
agent\_qwen3max\_vote\_vs\_selector & test & 268 & +4.9 [+1.9, +7.4] \\
agent\_qwen3max\_vote\_vs\_beam & test & 268 & +9.3 [+2.3, +14.6] \\
agent\_qwen3maxcritic\_s0\_vs\_selector & test & 268 & +4.1 [+1.1, +6.8] \\
agent\_qwen3maxcritic\_s0\_vs\_beam & test & 268 & +8.6 [+2.0, +13.6] \\
agent\_vs\_selector & test & 268 & +4.1 [+1.3, +7.6] \\
agent\_vs\_beam & test & 268 & +8.6 [+1.9, +13.4] \\
agent\_vs\_prev & test & 268 & +2.9 [-0.7, +5.3] \\
agent\_qwen3max\_vote\_vs\_selector & few-shot & 391 & +0.5 [-1.8, +3.4] \\
agent\_qwen3max\_vote\_vs\_beam & few-shot & 391 & -0.3 [-3.4, +4.5] \\
agent\_qwen3maxcritic\_s0\_vs\_selector & few-shot & 391 & +0.0 [-2.0, +2.1] \\
agent\_qwen3maxcritic\_s0\_vs\_beam & few-shot & 391 & -0.8 [-3.6, +3.0] \\
agent\_vs\_selector & few-shot & 391 & -0.7 [-2.8, +1.6] \\
agent\_vs\_beam & few-shot & 391 & -1.4 [-4.8, +3.1] \\
agent\_vs\_prev & few-shot & 391 & -2.6 [-4.8, +0.3] \\
agent\_vs\_selector & both folds & 659 & +1.3 [-0.7, +3.2] \\
agent\_vs\_beam & both folds & 659 & +2.6 [-1.8, +7.3] \\
agent\_vs\_prev & both folds & 659 & -0.4 [-2.8, +2.0] \\
\bottomrule\end{tabular}
\end{table}

\IfFileExists{tables/agent_selector_max_val.tex}{\begin{table}[t]\centering\scriptsize\setlength{\tabcolsep}{3pt}
\caption{Ceiling exploration: the selector over the ten presets and the extended diagnostics on the cached validation windows (cross-fitted by driver).}\label{tab:supp_agent_sel_max}
\begin{tabular}{lrrr}\toprule
row & windows & drivers & route \% [CI] \\\midrule
selector\_cv & 428 & 22 & 79 [74, 84] \\
beam1 & 428 & 22 & 77 [65, 83] \\
pf1 & 428 & 22 & 72 [65, 77] \\
ep1 & 428 & 22 & 24 [18, 41] \\
oracle & 428 & 22 & 92 [90, 94] \\
\bottomrule\end{tabular}
\end{table}
}{}
\IfFileExists{tables/agent_max_hints.tex}{\begin{table}[H]\centering\scriptsize\setlength{\tabcolsep}{3pt}
\caption{Calibration hints given to the C-max controller (validation candidates: share that recovered the route, by tool Top-1 and by diagnostic bin).}\label{tab:supp_agent_hints}
\begin{tabular}{llrr}\toprule
diagnostic & bin & candidates & route \% \\\midrule
tool\_top1 & beam & 428 & 77 \\
tool\_top1 & beam\_loose & 428 & 77 \\
tool\_top1 & beam\_wide & 428 & 77 \\
tool\_top1 & ep & 419 & 25 \\
tool\_top1 & ep\_hm & 426 & 22 \\
tool\_top1 & pf & 428 & 72 \\
tool\_top1 & pf\_s1 & 428 & 73 \\
tool\_top1 & pf\_s2 & 428 & 72 \\
tool\_top1 & pf\_s3 & 428 & 74 \\
tool\_top1 & pf\_s4 & 428 & 73 \\
consensus\_tools & 0 & 7{,}419 & 6 \\
consensus\_tools & 1 & 1{,}385 & 8 \\
consensus\_tools & 2-3 & 1{,}284 & 29 \\
consensus\_tools & 4-99 & 11{,}070 & 79 \\
hd\_abs\_rms\_deg & 0-10 & 400 & 99 \\
hd\_abs\_rms\_deg & 10-20 & 1{,}457 & 92 \\
hd\_abs\_rms\_deg & 20-40 & 4{,}779 & 74 \\
hd\_abs\_rms\_deg & >=40 & 14{,}817 & 31 \\
jn\_disagree & 0 & 7{,}006 & 56 \\
jn\_disagree & 1 & 5{,}276 & 46 \\
jn\_disagree & 2-99 & 9{,}171 & 38 \\
dr\_dist\_first60\_m & 0-20 & 4{,}130 & 51 \\
dr\_dist\_first60\_m & 20-50 & 6{,}845 & 53 \\
dr\_dist\_first60\_m & 50-150 & 6{,}631 & 52 \\
dr\_dist\_first60\_m & >=150 & 3{,}847 & 18 \\
turns\_unmatched & 0 & 4{,}732 & 86 \\
turns\_unmatched & 1 & 5{,}076 & 53 \\
turns\_unmatched & 2-99 & 11{,}645 & 27 \\
\bottomrule\end{tabular}
\end{table}
}{}}{}

\section{S6. Recovered windows}
\begin{figure*}[t]\centering
\includegraphics[width=0.9\textwidth]{fig_examples_w05.pdf}
\caption{Six held-out five-minute windows recovered by the offline beam and the online filter (speed, steering, start and
heading known) below their camera frames at mid-window (not an input); grey: roads; numbers: ADE.}
\end{figure*}
\begin{figure*}[t]\centering
\includegraphics[width=0.8\textwidth]{fig_method.pdf}
\caption{The pipeline on one held-out window: CAN signals, start, heading and visible roads in; candidate paths; scoring
at the reference seconds.}
\end{figure*}
\begin{figure*}[t]\centering
\includegraphics[width=0.75\textwidth]{fig_cdf_w05.pdf}
\caption{ADE distributions of the main rows on five-minute windows.}
\end{figure*}
\begin{figure*}[t]\centering
\includegraphics[width=0.75\textwidth]{fig_horizon_cohort.pdf}
\caption{Same-cohort horizon study: median ADE 1, 5 and 15 minutes after the start with driver-resampling bands.}
\end{figure*}

\bibliographystyle{IEEEtran}
\bibliography{../refs}
\end{document}
